% Standalone source generated from article.tex and its inputs.
% Edit the modular sources, then run python3 build.py.
% Begin preamble.tex
\documentclass[11pt,letterpaper]{article}
\usepackage[T1]{fontenc}
\usepackage{lmodern}
\usepackage[margin=0.92in,headheight=15pt]{geometry}
\usepackage{amsmath,amssymb,amsthm,mathtools,mathrsfs}
\usepackage{booktabs,longtable,array,enumitem,xcolor,microtype,xurl}
\usepackage{fancyhdr,needspace}
\usepackage[colorlinks=true,linkcolor=blue!40!black,
 citecolor=teal!55!black,urlcolor=blue!50!black]{hyperref}
\usepackage{bookmark}
\definecolor{ink}{RGB}{25,47,68}
\emergencystretch=2em
\setlength{\parskip}{3pt}
\setlength{\parindent}{1.25em}
\setcounter{tocdepth}{2}
\numberwithin{equation}{section}
\theoremstyle{plain}
\newtheorem{theorem}{Theorem}[section]
\newtheorem{lemma}[theorem]{Lemma}
\newtheorem{proposition}[theorem]{Proposition}
\newtheorem{corollary}[theorem]{Corollary}
\theoremstyle{definition}
\newtheorem{definition}[theorem]{Definition}
\newtheorem{example}[theorem]{Example}
\theoremstyle{remark}
\newtheorem{remark}[theorem]{Remark}
\DeclareMathOperator{\Li}{Li}
\DeclareMathOperator{\FP}{FP}
\DeclareMathOperator{\CT}{CT}
\DeclareMathOperator{\Log}{Log}
\DeclareMathOperator{\Res}{Res}
\DeclareMathOperator{\spanQ}{span_{\mathbb Q}}
\newcommand{\dd}{\,\mathrm d}
\newcommand{\C}{\mathbb C}
\newcommand{\Q}{\mathbb Q}
\newcommand{\N}{\mathbb N}
\newcommand{\Z}{\mathbb Z}
\newcommand{\HH}{\mathcal H}
\newcommand{\src}[1]{\nolinkurl{#1}}
\newcommand{\st}{\boldsymbol s}
\newcommand{\tz}{\boldsymbol t}
\newcommand{\xz}{\boldsymbol x}
\newcommand{\uz}{\boldsymbol u}
\pagestyle{fancy}
\fancyhf{}
\fancyhead[L]{\small\textsc{Harmonic gaps and mixed moments}}
\fancyhead[R]{\small ProveIt research continuation}
\fancyfoot[C]{\thepage}
\renewcommand{\headrulewidth}{0.3pt}
\hypersetup{pdftitle={Harmonic Gaps, Gamma Derivatives, and Mixed Tail Moments},
 pdfauthor={ProveIt research continuation},
 pdfsubject={Joint Lerch generators, transverse harmonic derivatives, and exact centered moments}}


% End preamble.tex

\begin{document}
\hypersetup{pageanchor=false}
\begin{titlepage}
\thispagestyle{empty}
{\color{ink}\small\textsc{ProveIt \quad / \quad Research continuation}\par}
\vspace{1.1cm}
{\color{ink}\huge\bfseries Harmonic Gaps, Gamma Derivatives,\\[5pt]
and Mixed Tail Moments\par}
\vspace{0.45cm}
{\Large Joint Lerch generators and exact centered values\par}
\vspace{0.8cm}
{\large 11 October 2026\par}
\vspace{0.5cm}
\noindent\textbf{Repository baseline.}\quad
\src{7bd45777a8a127e5dc69aff8887ca36a79a3059d}.
The manuscript and all eleven archives in the incoming directory were
inspected at this revision.
\vspace{0.7cm}

\begin{abstract}
We address three connected questions in the newest harmonic-parity report.
First, arbitrary zero blocks around any finite number of independently
moving harmonic orders admit one jointly entire multiple-Lerch completion.
A confluent divided-difference kernel gives a common tube of holomorphy;
the complete decoration series converges on the unit polydisc, independently
of the number of free orders. Second, a derivative transverse to a
nonpositive harmonic index is represented by an explicit convergent
logarithmic Mellin integral and by a finite difference of decorated
lower-depth sums. At a zero index the first new gap is a log-Gamma
difference; higher derivatives give normalized Stieltjes primitives.
All derivative orders, all nonpositive centers, and arbitrary unmarked
nonpositive decorations are included. Third, every centered even moment
of a product of two even-order Hurwitz tails is evaluated by a finite
Bernoulli formula. Classical odd-weight Euler reduction puts the values
in an explicit finite span of ordinary zeta products. One analytic
generating function produces the entire moment family and has finite
polylogarithmic primitives.

The underlying harmonic interpolant, its untwisted geometric kernel, and
the classical Euler and Hurwitz identities are credited to their sources.
The contribution is the stated extension of the repository's exact
questions, not a claim of literature-wide priority. Proofs are separated
from exact finite checks and high-precision diagnostics. The Gaussian
$S_6$ and revised $S_8$ evaluations remain conjectural.
\end{abstract}
\vfill
\noindent\small
\textbf{Contents of the accompanying package.}
Modular and standalone \LaTeX\ sources, this PDF, exact verification
programs, numerical diagnostics, result records, source provenance,
an integration ledger, and further research questions.
\end{titlepage}
\pagenumbering{roman}
\hypersetup{pageanchor=true}
\tableofcontents
\clearpage
\pagenumbering{arabic}
% Begin sections/01_scope.tex
\section{Scope, principal results, and conventions}
\label{scope:section}

This continuation starts at the explicit questions R1, R7, and R8 of
\emph{Harmonic Parity and Resolvent Identities}, preserved in the incoming
directory at the pinned revision \cite{IncomingParity,Repository}.
R1 asks for mixed even-tail moments, beginning with orders two and four.
R7 asks for a genuinely convergent generator with several free harmonic
orders and arbitrary zero blocks. R8 asks for a derivative transverse to
a harmonic index that is subsequently set to a nonpositive integer.
The corresponding results here are independent theorems with compatible
normalizations.

\subsection{What is established}

\begin{enumerate}[leftmargin=*,label=\textbf{\arabic*.}]
\item \textbf{All zero decorations with independent orders.}
Theorem~\ref{zbg:main} proves a single analytic formula for
\[
 \sum_{p_0,\ldots,p_k\geq0}
 \HH_{a,b}(0^{p_0},s_1,0^{p_1},\ldots,s_k,0^{p_k})
                     u_0^{p_0}\cdots u_k^{p_k}.
\]
Every $s_j$ is an independent complex variable. The series converges
locally normally for $|u_j|<1$, with a domain that does not shrink with
$k$. The proof gives all branch cancellations and does not infer an
analytic identity from agreement at integer endpoints.

\item \textbf{Transverse Gamma and Stieltjes gaps.}
Theorems~\ref{tr:gap} and~\ref{tr:integral} prove every order of
\[
 \left.\partial_u^r\HH_{a,b}(s,u,t)\right|_{u=-m},
 \qquad m\geq0,\quad r\geq1.
\]
The first new case is the log-Gamma gap
$\log\Gamma(z+1)-\log\Gamma(y)$ between two surviving lattice variables
$y>z$. There are explicit absolutely convergent Mellin integrals, finite
decorated-sum formulas, and a finite family sufficient after arbitrary
unmarked nonpositive indices have been removed. The simultaneous
Stieltjes completion at $u=1$ is included.

\item \textbf{Every mixed even-tail moment.}
Section~\ref{mt:section} evaluates all centered even moments of
$(\zeta(A)-H_n^{(A)})(\zeta(B)-H_n^{(B)})$ for even $A,B\geq2$.
The finite formula uses Bernoulli numbers and classical double zeta
values of odd weight; Euler's reduction removes the latter.
For each fixed pair $(A,B)$ a fixed finite span contains the entire
infinite sequence of moments. An analytic generating function and
finite ordinary-polylogarithm primitives complete the result.
This resolves the two-factor sector of R1; products with additional
harmonic factors are not covered by the same theorem.
\end{enumerate}

For a concrete illustration, put $x_n=n+\tfrac12$ and
$T_j(n)=\zeta(j,n+1)$. The unequal-order identities include
\begin{align}
 \sum_{n\geq0}T_2(n)T_4(n)
  &=\frac{15}{2}\zeta(5)-3\zeta(2)\zeta(3),\label{scope:first}\\
 \sum_{n\geq0}x_n^2T_2(n)T_4(n)
  &=\frac14\zeta(2)\zeta(3)+\frac16\zeta(2)
       +\frac12\zeta(3)-\frac58\zeta(5),\label{scope:second}\\
 \sum_{n\geq0}\left\{x_n^4T_2(n)T_4(n)-\frac13\right\}
  &=-\frac7{80}\zeta(2)\zeta(3)-\frac1{30}\zeta(2)
       -\frac14\zeta(3)+\frac7{32}\zeta(5)-\frac3{40}.
       \label{scope:third}
\end{align}
All three displayed series, including the brackets in the last one,
are ordinary convergent series. The third identity must retain its
subtraction. Since $T_2(n)=\psi'(n+1)$ and
$T_4(n)=\psi'''(n+1)/6$, multiplication by six gives corresponding
polygamma-product evaluations.

\subsection{Provenance and the precise claim of advancement}

The harmonic interpolant and its original geometric Mellin kernel are
established mathematics. Ihara, Nakamura, and Yamamoto prove their
entire interpolation and harmonic product law; their paper also
explains related earlier work of Komori
\cite[Proposition 3.3, Theorem 4.1, Corollary 4.2]{INY}.
The incoming independent-orders report identifies the same interpolant
and gives a relative endpoint normalization \cite{IncomingIndependent}.
The zero-index polynomial reductions and one-free-order generator are
already in \cite{IncomingParity}. We retain those definitions and give
the requested simultaneous analytic extension and transverse formulas.

Similarly, Bernoulli power sums, Hurwitz specializations, Gamma
asymptotics, and Euler's odd-weight double-zeta reduction are classical
\cite{DLMFHurwitz,DLMFGamma,DLMFBernoulli,Bradley2007}.
They are ingredients of the mixed-tail theorem, rather than separate
claims of new mathematics.

The present results are proposed additions relative to the inspected
repository. We make no assertion of literature-wide priority for every
specialization, no arithmetic-independence claim for a spanning set, and
no proof-assistant formalization claim. The original $S_6$ and revised
$S_8$ candidates are not proved or refuted here.

\subsection{Analytic and algebraic conventions}

All endpoint parameters $a,b$ lie in the right half-plane.
Principal logarithms define $(n+a)^{-s}$. Positive real integration
variables use their real logarithms. The Hurwitz zeta function is
normalized by
\begin{equation}\label{scope:stieltjes}
 \zeta(1+v,x)=\frac1v+
       \sum_{r=0}^{\infty}\frac{(-1)^r}{r!}\gamma_r(x)v^r,
 \qquad \gamma_0(x)=-\psi(x).
\end{equation}
A superscript $\zeta^{(r)}$ denotes differentiation in the first
argument. In particular \cite[\S25.11]{DLMFHurwitz},
\begin{equation}\label{scope:hurwitz}
 \zeta(-m,x)=-\frac{B_{m+1}(x)}{m+1},\qquad
 \zeta'(0,x)=\log\Gamma(x)-\frac12\log(2\pi).
\end{equation}
Bernoulli numbers mean $B_j=B_j(0)$, so $B_1=-1/2$ and $B_1(1)=1/2$.
The strict multiple-zeta convention is
\[
 \zeta(r,t)=\sum_{n>j\geq1}n^{-r}j^{-t}.
\]

For a word $w=(s_1,\ldots,s_k)$, define the meromorphic ordered tail
\[
 Z_c(w)=\sum_{n_1>\cdots>n_k\geq0}
                      \prod_j(n_j+c)^{-s_j},\qquad Z_c(\varnothing)=1.
\]
An initial domain with $\Re s_j>1$ suffices below. The relative
interpolant is the unique triangular solution of
\begin{equation}\label{scope:relative}
 Z_b(w)=\sum_{w=uv}Z_a(u)\HH_{a,b}(v),\qquad
 \HH_{a,b}(\varnothing)=1.
\end{equation}
The sum includes both empty prefix and empty suffix. Its depth-one value is
\begin{equation}\label{scope:h}
 h_{a,b}(s)=\zeta(s,b)-\zeta(s,a).
\end{equation}
The pole at $s=1$ cancels. Useful values are
\begin{align}
 h_{a,b}(1)&=\psi(a)-\psi(b),&
 h_{a,b}^{(r)}(1)&=(-1)^r\{\gamma_r(b)-\gamma_r(a)\},\label{scope:hjets}\\
 h_{a,b}'(0)&=\log\Gamma(b)-\log\Gamma(a),&
 h_{a,b}(0)&=a-b.\notag
\end{align}
At $a=b+N$, $N\in\N_0$, the relative word is the finite strict sum
over $N>n_1>\cdots>n_k\geq0$. For general endpoints it means
\eqref{scope:relative}, with the joint continuation proved again in
Section~\ref{zbg:section}. The finite endpoint property alone is not
our definition of continuation.

We use a cutoff constant only for a function with a specified asymptotic
expansion in powers and logarithms of that cutoff. It is the coefficient
of the monomial $N^0(\log N)^0$. Section~\ref{mt:section} proves when
that prescription agrees with the spectral value. This distinction
is essential: even an accidentally convergent paired boundary series
need not equal the analytic continuation from its initial half-plane.


% End sections/01_scope.tex

% Begin sections/02_joint_generator.tex
\section{A jointly entire kernel for all zero-block decorations}
\label{zbg:section}

The one-free-order zero-block generator in the incoming harmonic-parity
report admits a uniform extension to any number of independently moving
orders.  The extension is not obtained by interpolating finite sums from
integer endpoint differences.  A divided-difference identity proves the
required analytic relation at arbitrary endpoints before any coefficient
is extracted.

The underlying geometric kernel and its normalized Mellin interpolation
are already present in Ihara, Nakamura, and Yamamoto
\cite[Proposition~3.3 and Theorem~4.1]{INY}; the published paper is
\emph{The Ramanujan Journal} \textbf{67} (2025), article~1,
doi:10.1007/s11139-025-01045-2.  In their notation the kernel below is
$e^{-b\sum_jx_j}G_k(e^{-x_1},\ldots,e^{-x_k};a-b)$.
Equivalently it is $e^{-(b-1)\sum_jx_j}$ times their $b=1$ kernel.
Its divided-difference form, composition proof, and continuation proof
are included to make the normalization and later estimates explicit;
the underlying interpolant and its geometric Mellin representation are
not claimed as new.  The developments here are the simultaneous twist
tube, the full dimension-independent zero-decoration polydisc, the
all-block identity with arbitrary independent free orders, and its
compatibility with the transverse derivative calculus developed below.

Write $\mathcal H_{a,b}(s_1,\ldots,s_k)$ for the entire relative
Hurwitz interpolant, normalized by
\[
 Z_b(w)=\sum_{w=uv} Z_a(u)\mathcal H_{a,b}(v),\qquad
 Z_c(s_1,\ldots,s_j)=\sum_{n_1>\cdots>n_j\geq0}
                         \prod_{\ell=1}^j(n_\ell+c)^{-s_\ell}.
\]
Empty words have value one.  Throughout this section
$\Re a,\Re b>0$, and powers of the positive-real-axis variables use
their real logarithms.  All endpoint powers use the branches obtained
on the right half-plane.

\subsection{The established geometric kernel in divided-difference form}

For $\boldsymbol x=(x_1,\ldots,x_k)$ put
\[
 X_0=0,\qquad X_j=x_1+\cdots+x_j,\qquad Q_j=e^{-X_j},\qquad
 \Delta=a-b.
\]
At pairwise distinct nodes define
\begin{equation}\label{zbg:kernel}
 \mathscr K_{a,b}^{(k)}(\boldsymbol x)
 =e^{-bX_k-\sum_{j=1}^{k-1}X_j}
       \sum_{j=0}^k
       \frac{e^{-\Delta X_j}}{\prod_{\ell\ne j}(Q_j-Q_\ell)}.
\end{equation}
The sum is the divided difference
$[Q_0,\ldots,Q_k]z^\Delta$, where the value at $Q_j$ is specified
as $e^{-\Delta X_j}$.  This specification keeps track of the
logarithms even when the exponential nodes wind around zero.  Set
$\mathscr K^{(0)}=1$.

\begin{lemma}[Kernel analyticity and exponential control]\label{zbg:bounds}
The expression \eqref{zbg:kernel} extends holomorphically through
every hyperplane $X_j-X_i=0$, away from the nonzero-period
hyperplanes below. Its only possible finite singular
hyperplanes are
\[
 X_j-X_i=2\pi i\ell,\qquad 0\leq i<j\leq k,
                   \quad \ell\in\mathbb Z\setminus\{0\}.
\]
Define the open convex tube
\begin{equation}\label{zbg:tube}
 \Omega_k=\left\{\boldsymbol t\in\mathbb C^k:
   \left|\Im\sum_{j=p}^q t_j\right|<2\pi
                      \quad(1\leq p\leq q\leq k)\right\}.
\end{equation}
On compact subsets of $\Re a,\Re b>0$ and
$\boldsymbol t\in\Omega_k$, for every multi-index $\boldsymbol\nu$
there are $C,\delta>0$ such that
\begin{equation}\label{zbg:bound}
 \left|\partial_{\boldsymbol x}^{\boldsymbol\nu}
       \mathscr K_{a,b}^{(k)}(\boldsymbol x-\boldsymbol t)\right|
       \leq C e^{-\delta(x_1+\cdots+x_k)},\qquad x_j\geq0.
\end{equation}
\end{lemma}

\begin{proof}
If nodes collide with $X_j-X_i=0$, their assigned logarithms agree.
On a sufficiently small neighborhood of each cluster, choose the common
holomorphic logarithm.  The ordinary confluent divided difference of
its holomorphic power function extends across all coinciding nodes.
Different clusters have nonzero denominator differences.  This proves
the simultaneous removability assertion, including intersections of
arbitrarily many zero hyperplanes away from the remaining singular
hyperplanes. Nonzero exponential periods are
exactly the remaining possible denominator zeros.

The first-order kernel is
\[
 \mathscr K_{a,b}^{(1)}(v)=\frac{e^{-bv}-e^{-av}}{1-e^{-v}}.
\]
The elementary divided-difference recurrence gives, for $k\geq2$,
\begin{equation}\label{zbg:recurrence}
 \mathscr K_{a,b}^{(k)}(v,w,\boldsymbol z)=
 \frac{e^{-v}\mathscr K_{a,b}^{(k-1)}(v+w,\boldsymbol z)
       -e^{-av}\mathscr K_{a,b}^{(k-1)}(w,\boldsymbol z)}
      {1-e^{-v}}.
\end{equation}
Indeed the two divided differences in the numerator omit $Q_1$ and
$Q_0$, respectively; scaling the nodes of the second one by $Q_1$
gives exactly the displayed exponential factor.

This recurrence proves \eqref{zbg:bound} by induction.  For large
$x_1$, the denominator is uniformly bounded away from zero, and both
numerator terms have the required exponential decay.  For bounded
$x_1$, the only possible denominator zero is the removable zero at
$v=0$; the numerator vanishes there because its two kernel terms agree.  A fixed small neighborhood of the shifted positive orthant
stays away from all nonzero-period hyperplanes, because the compact
subset in \eqref{zbg:tube} has a positive distance from its boundary.
The maximum principle, or the Cauchy formula around the removable
zero, then bounds the quotient by the inductively bounded numerator.
Merging $v+w$ preserves the tube inequalities, since they involve
consecutive sums.  One may take $\delta$ smaller than the minimum of
$1,\Re a,\Re b$ on the endpoint compact set.  Cauchy estimates in the
same uniform neighborhood prove the bounds for every derivative.
\end{proof}

\subsection{Why the kernel gives the specified noninteger interpolant}

Let
\begin{equation}\label{zbg:infinite}
 \mathscr B_c^{(k)}(\boldsymbol x)
 =\frac{e^{-cX_k-\sum_{j=1}^{k-1}X_j}}
              {\prod_{j=1}^k(1-e^{-X_j})},\qquad x_j>0.
\end{equation}
This is the geometric kernel of the ordered infinite sum $Z_c$:
introducing the positive successive gaps gives
$\mathscr B_c^{(k)}=\sum_{n_1>\cdots>n_k\geq0}
e^{-\sum_j(n_j+c)x_j}$.

\begin{lemma}[Exact endpoint composition]\label{zbg:composition}
For arbitrary right-half-plane endpoints $a,b,c$,
\begin{equation}\label{zbg:kernelcomposition}
 \mathscr K_{a,b}^{(k)}(x_1,\ldots,x_k)
 =\sum_{j=0}^k
   \mathscr K_{a,c}^{(j)}(x_1,\ldots,x_j)
   \mathscr K_{c,b}^{(k-j)}(x_{j+1},\ldots,x_k).
\end{equation}
Consequently,
\begin{equation}\label{zbg:convolution}
 \mathscr B_b^{(k)}=\sum_{j=0}^k
                       \mathscr B_a^{(j)}\mathscr K_{a,b}^{(k-j)},
\end{equation}
with the same prefix/suffix arguments.
\end{lemma}

\begin{proof}
Apply the divided-difference product identity
\[
 [Q_0,\ldots,Q_k](fg)=\sum_{j=0}^k
 [Q_0,\ldots,Q_j]f\,[Q_j,\ldots,Q_k]g
\]
to $f(z)=z^{a-c}$ and $g(z)=z^{c-b}$.  The suffix nodes in its
kernel are $Q_\ell/Q_j$.  Scaling a divided difference of order
$k-j$ supplies the factor $Q_j^{k-j-(c-b)}$; together with the
kernel prefactors this makes every summand have the common factor
$e^{-bX_k-\sum_{\ell=1}^{k-1}X_\ell}$.  Thus the product identity
is exactly \eqref{zbg:kernelcomposition}.  This argument is first
made at distinct positive nodes and then continued through the
removable collisions.

For $x_j>0$, let a real upper endpoint $A$ tend to $+\infty$ in
$\mathscr K_{A,b}=\mathscr K_{A,a}*\mathscr K_{a,b}$.  In
\eqref{zbg:kernel}, every term except $j=0$ tends to zero.  The
remaining term is precisely \eqref{zbg:infinite}.  This gives
\eqref{zbg:convolution}, a finite identity requiring no interchange
with an infinite depth sum.  In particular, the proof is valid for
noninteger $a-b$ and does not use uniqueness of interpolation from
integer values.
\end{proof}

\subsection{The entire multiple-Lerch completion}

For $\Re s_j>0$ define
\begin{equation}\label{zbg:Mellin}
 \mathscr E_{a,b}(\boldsymbol s;\boldsymbol t)
 =\frac1{\prod_{j=1}^k\Gamma(s_j)}
   \int_{(0,\infty)^k}\prod_{j=1}^k x_j^{s_j-1}
      \mathscr K_{a,b}^{(k)}(\boldsymbol x-\boldsymbol t)
           \,d\boldsymbol x,\qquad\boldsymbol t\in\Omega_k.
\end{equation}

\begin{theorem}[Independent spectral orders and simultaneous branch cancellation]
\label{zbg:entire}
The function \eqref{zbg:Mellin} extends jointly to an entire function
of $\boldsymbol s\in\mathbb C^k$, holomorphic in
$\boldsymbol t\in\Omega_k$, with locally normal derivatives of
every fixed spectral and decoration order.  It satisfies
\begin{equation}\label{zbg:lowering}
 \mathscr E_{a,b}(\boldsymbol s;\boldsymbol0)
       =\mathcal H_{a,b}(\boldsymbol s),\qquad
 \partial_{t_j}\mathscr E_{a,b}(\boldsymbol s;\boldsymbol t)
       =\mathscr E_{a,b}(\boldsymbol s-\boldsymbol e_j;\boldsymbol t).
\end{equation}
In a neighborhood of $\boldsymbol t=0$ its normally convergent Taylor
series is
\begin{equation}\label{zbg:Taylor}
 \mathscr E_{a,b}(\boldsymbol s;\boldsymbol t)
 =\sum_{\boldsymbol\nu\geq0}
     \mathcal H_{a,b}(\boldsymbol s-\boldsymbol\nu)
                          \frac{\boldsymbol t^{\boldsymbol\nu}}{\boldsymbol\nu!}.
\end{equation}
For instance, the Taylor series converges normally whenever
$\sum_j|t_j|<2\pi$.
\end{theorem}

\begin{proof}
Lemma~\ref{zbg:bounds} proves absolute convergence of the integral.
Repeated integration by parts gives the particularly useful continuation
formula
\begin{equation}\label{zbg:continuation}
 \mathscr E_{a,b}(\boldsymbol s;\boldsymbol t)
 =\frac{(-1)^{|\boldsymbol N|}}
             {\prod_j\Gamma(s_j+N_j)}
    \int_{(0,\infty)^k}\prod_jx_j^{s_j+N_j-1}
      \partial_{\boldsymbol x}^{\boldsymbol N}
          \mathscr K_{a,b}^{(k)}(\boldsymbol x-\boldsymbol t)
                                                    \,d\boldsymbol x,
\end{equation}
valid for $\Re s_j>-N_j$.  It agrees with the original integral
where $\Re s_j>0$; the endpoint terms there vanish.  Since each
$N_j$ is arbitrary, these compatible formulas prove joint entireness.
The same exponential majorants, with logarithmic powers from spectral
differentiation, prove local normality.

At $\boldsymbol t=0$, apply the normalized Mellin transform to
\eqref{zbg:convolution}, initially with all spectral real parts
sufficiently large.  The product of prefix/suffix transforms is a
product of independent integrals.  The result is exactly the defining
deconcatenation relation for $\mathcal H_{a,b}$.  This proves the
first identity in \eqref{zbg:lowering}; entireness extends it everywhere.
The equality $\partial_{t_j}\mathscr K=-\partial_{x_j}\mathscr K$
and one integration by parts prove the lowering relation first where
$\Re s_j>1$, and then everywhere.  Taylor's theorem proves
\eqref{zbg:Taylor}.  Every polydisc whose coordinate radii have sum
less than $2\pi$ lies in $\Omega_k$, giving the stated domain.
\end{proof}

For $\boldsymbol t\in\Omega_k$ with $\Re t_j<0$ for every $j$, let
\[
 \mathscr Z_c(\boldsymbol s;\boldsymbol t)
 =\sum_{n_1>\cdots>n_k\geq0}
       \prod_j(n_j+c)^{-s_j}e^{(n_j+c)t_j}.
\]
These ordered multiple-Lerch series (with the standard depth-one
normalization \cite[\S25.14]{DLMFLerch}) are absolutely convergent and
entire in their independent orders.  In that domain,
$\mathscr E_{a,b}=\mathscr Z_a^{-1}*\mathscr Z_b$ under
deconcatenation, by the same kernel identity.  At each fixed depth the
inverse is a finite expression.  For example,
\begin{align*}
 \mathscr E_{a,b}(s,t;u,v)={}&
  \mathscr Z_b(s,t;u,v)-\mathscr Z_a(s,t;u,v)\\
 &-\mathscr Z_a(s;u)
               \bigl(\mathscr Z_b(t;v)-\mathscr Z_a(t;v)\bigr).
\end{align*}
Theorem~\ref{zbg:entire} proves that the whole expression has a
single holomorphic continuation throughout $\Omega_k$, including all
simultaneous zero-twist branch crossings.  Continuing the summands
independently without this completion is unnecessary.

\subsection{Summing every zero block on a dimension-independent polydisc}

Let $\boldsymbol u=(u_0,\ldots,u_k)$, $|u_j|<1$, and take the
logarithms $L_j=\log(1+u_j)$ determined at the origin.  Put
$t_j=L_j-L_{j-1}$.

\begin{theorem}[All zero decorations around independent orders]
\label{zbg:main}
For every $k\geq1$,
\begin{align}\label{zbg:generator}
 &\sum_{p_0,\ldots,p_k\geq0}
 \mathcal H_{a,b}(\{0\}^{p_0},s_1,\{0\}^{p_1},\ldots,
                       s_k,\{0\}^{p_k})\prod_{j=0}^ku_j^{p_j}
 \notag\\
 &\qquad=(1+u_0)^{a-1}(1+u_k)^{-b}
                    \prod_{j=1}^{k-1}(1+u_j)^{-1}
        \mathscr E_{a,b}(\boldsymbol s;
                L_1-L_0,\ldots,L_k-L_{k-1}).
\end{align}
The series converges absolutely and locally normally on the full
polydisc $|u_j|<1$, uniformly on compact spectral and endpoint sets.
Every fixed collection of independent spectral derivatives may be
taken term by term.
\end{theorem}

\begin{proof}
For $|u|<1$, $1+u$ lies in the right half-plane, so
$|\Im\log(1+u)|<\pi/2$.  Every consecutive sum of the $t_j$
telescopes to $L_q-L_{p-1}$ and has imaginary part of absolute value
less than $\pi$.  Thus the right side of \eqref{zbg:generator}
is holomorphic on the entire stated polydisc, independently of $k$.

To identify its Taylor coefficients, use \eqref{zbg:Taylor} locally
and formally replace
$\mathcal H_{a,b}(\boldsymbol s-\boldsymbol\nu)$ by
$\boldsymbol y^{\boldsymbol\nu}$.  The resulting expression is
\[
 (1+u_0)^{a-y_1-1}
 \prod_{j=1}^{k-1}(1+u_j)^{y_j-y_{j+1}-1}
 (1+u_k)^{y_k-b}.
\]
Its coefficient of $\prod u_j^{p_j}$ is
\[
 \binom{a-y_1-1}{p_0}
 \prod_{j=1}^{k-1}\binom{y_j-y_{j+1}-1}{p_j}
 \binom{y_k-b}{p_k}.
\]
This is exactly the canonical strict-gap polynomial of the zero
blocks.  The analytic gap-deletion theorem in \cite{IncomingParity},
with the independent-order continuation fixed in \cite{IncomingIndependent},
identifies its finite polynomial expansion with
the left coefficient in \eqref{zbg:generator}.  Each coefficient
uses only finitely many terms of \eqref{zbg:Taylor}; no interchange
of a divergent formal order-shift sum is involved.  Holomorphy of the
right side now proves absolute local convergence of its Taylor series
throughout the unit polydisc.  Spectral derivatives obey the same
Cauchy estimates uniformly on compact spectral sets.
\end{proof}

\begin{corollary}[Insertion sums and normalized primitives]\label{zbg:corollaries}
For every integer $r\geq0$,
\begin{equation}\label{zbg:insertion}
 \sum_{p_0+\cdots+p_k=r}
 \mathcal H_{a,b}(\{0\}^{p_0},s_1,\ldots,s_k,\{0\}^{p_k})
 =\binom{a-b-k}{r}\mathcal H_{a,b}(s_1,\ldots,s_k),
\end{equation}
where the zero blocks between consecutive displayed free orders are
included in the summand as in \eqref{zbg:generator}.  Moreover,
\begin{equation}\label{zbg:primitive}
 \int_{\tau_0}^{t_j}\mathscr E_{a,b}(\boldsymbol s;
                      t_1,\ldots,\tau,\ldots,t_k)\,d\tau
 =\mathscr E_{a,b}(\boldsymbol s+\boldsymbol e_j;\boldsymbol t)
  -\mathscr E_{a,b}(\boldsymbol s+\boldsymbol e_j;
                     t_1,\ldots,\tau_0,\ldots,t_k),
\end{equation}
along any path in the indicated slice of $\Omega_k$.
For $\boldsymbol m\in\mathbb N_0^k$,
\begin{equation}\label{zbg:negative}
 \mathscr E_{a,b}(-\boldsymbol m;\boldsymbol t)
 =(-1)^{|\boldsymbol m|}
       \partial_{\boldsymbol x}^{\boldsymbol m}
            \mathscr K_{a,b}^{(k)}(-\boldsymbol t).
\end{equation}
\end{corollary}

\begin{proof}
Set every $u_j=u$ in \eqref{zbg:generator}.  All twists vanish and
the prefactor becomes $(1+u)^{a-b-k}$.  The primitive is
\eqref{zbg:lowering} with the additive normalization retained.
For \eqref{zbg:negative}, use \eqref{zbg:continuation} with
$N_j=m_j+1$ and integrate the derivatives successively.  Exponential
decay removes the upper boundary terms and supplies precisely one
additional minus sign for each integration.
\end{proof}

At $a=b+N$ with $N\geq0$ integral, the kernel is the finite ordered
exponential sum and $\mathscr E_{a,b}$ is its weighted finite
Dirichlet sum.  The generating function is consequently a polynomial
of total degree at most $N-k$ if $N\geq k$, and is zero if $N<k$.
For $k=1$ the theorem specializes to the previously proved
single-Lerch generator.  For $k>1$ it resolves the requested joint
convergence and simultaneous branch-cancellation problem.  It supplies
an explicit Gamma-separated kernel, with no claim that this kernel is
minimal in an unspecified class.

% End sections/02_joint_generator.tex

% Begin sections/03_transverse.tex
\section{Transverse index derivatives and Gamma--Stieltjes gaps}
\label{tr:section}

A formula proved only after setting an index to $-m$ cannot be
differentiated in that index. The derivative sees information that a
Bernoulli polynomial does not contain. This section constructs that
information explicitly. We first obtain identities valid with an
independent moving index, and only then differentiate and specialize.
The kernel of Section~\ref{zbg:section} supplies a second, convergent
integral realization of the same derivatives.

\subsection{The joint Hurwitz completion}

Put
\begin{equation}\label{tr:Q}
 Q(u,x)=\zeta(u,x)-\frac1{u-1},
\end{equation}
with its removable value at $u=1$. This is entire in $u$ and holomorphic
for $\Re x>0$. At the resonant point,
\begin{equation}\label{tr:Qstieltjes}
 \left.\partial_u^rQ(u,x)\right|_{u=1}=(-1)^r\gamma_r(x).
\end{equation}
For $m,r\geq0$ set
\begin{equation}\label{tr:ell}
 \ell_{m,r}(x)=\left.\partial_u^r\zeta(u,x)\right|_{u=-m}.
\end{equation}
The derivatives of $Q$ at $-m$ differ from $\ell_{m,r}$ only by the
constant $r!/(m+1)^{r+1}$. Thus they give the same gap differences.

For $\Re s$ sufficiently large, uniformly on compact sets of $u$,
define the paired sum
\begin{equation}\label{tr:Dtilde}
 \widetilde D_{a,b}(s,u)=
 \sum_{n\geq0}\left\{
 (n+b)^{-s}Q(u,n+b)-(n+a)^{-s}Q(u,n+a)\right\}.
\end{equation}
We initially require convergence of the separate one-sided series,
not merely an accidental cancellation at an isolated exponent.

\begin{theorem}[A jointly entire weighted-tail difference]
\label{tr:jointD}
The continuation of \eqref{tr:Dtilde} is the entire function
\begin{equation}\label{tr:Djoint}
 \boxed{\widetilde D_{a,b}(s,u)
       =Q(u,b)h_{a,b}(s)-\HH_{a,b}(s,u).}
\end{equation}
Consequently the series
\begin{equation}\label{tr:D}
 D_{m,r}(s;a,b)=\sum_{n\geq0}\left\{
 (n+b)^{-s}\ell_{m,r}(n+b)
 -(n+a)^{-s}\ell_{m,r}(n+a)\right\},
 \qquad \Re s>m+2,
\end{equation}
has the entire continuation
\begin{equation}\label{tr:Didentity}
 D_{m,r}(s;a,b)=\ell_{m,r}(b)h_{a,b}(s)
       -\left.\partial_u^r\HH_{a,b}(s,u)\right|_{u=-m}.
\end{equation}
For each $r\geq0$, the similarly continued Stieltjes-weighted difference is
\begin{align}
 &\sum_{n\geq0}\left\{
 \frac{\gamma_r(n+b)}{(n+b)^s}
 -\frac{\gamma_r(n+a)}{(n+a)^s}\right\}
 \notag\\
 &\qquad=\gamma_r(b)h_{a,b}(s)
       -(-1)^r\left.\partial_u^r\HH_{a,b}(s,u)\right|_{u=1}.
 \label{tr:StieltjesD}
\end{align}
The last series is first interpreted, for example, on $\Re s>1$.
\end{theorem}

\begin{proof}
In a domain of absolute convergence, summing the inner index gives
\[
 Z_c(s,u)=\zeta(s,c)\zeta(u,c)
          -\sum_{n\geq0}(n+c)^{-s}\zeta(u,n+c).
\]
Insert this identity into
\[
 \HH_{a,b}(s,u)=Z_b(s,u)-Z_a(s,u)-\zeta(s,a)h_{a,b}(u).
\]
The endpoint products cancel to leave
\[
 \sum_{n\geq0}\left\{(n+b)^{-s}\zeta(u,n+b)
               -(n+a)^{-s}\zeta(u,n+a)\right\}
 =\zeta(u,b)h_{a,b}(s)-\HH_{a,b}(s,u).
\]
Subtract $h_{a,b}(s)/(u-1)$ to obtain \eqref{tr:Djoint}.
Its right side is jointly entire by Theorem~\ref{zbg:entire}.

The Hurwitz asymptotic expansion, uniformly differentiated on a compact
neighborhood of $u=-m$, gives
$\ell_{m,r}(n+c)=O(n^{m+1}(1+\log^r n))$ on compact endpoint sets.
A slightly enlarged compact set gives termwise spectral differentiation
in a sufficiently far right half-plane; continuation then proves
\eqref{tr:Didentity}. Formula \eqref{tr:Qstieltjes} gives
\eqref{tr:StieltjesD}. At $u=1$, the coefficient $\gamma_r(n+c)$ is
$O(1+\log^{r+1}n)$ on compact endpoint sets, proving the stated
initial half-plane.
\end{proof}

The whole function \eqref{tr:Djoint}, rather than two separately
regularized tails, fixes all constants. In particular the theorem covers
simultaneous spectral derivatives at $s=1,u=1$.

\subsection{A finite gap identity before specialization}

For lattice functions $f,g$ with enough decay, define
\begin{align}
 Z_c[f]&=\sum_{n\geq0}f(n+c),&
 Z_c[f,g]&=\sum_{n>j\geq0}f(n+c)g(j+c),\notag\\
 \HH_{a,b}[f,g]
   &=Z_b[f,g]-Z_a[f,g]-Z_a[f]\{Z_b[g]-Z_a[g]\}.
 \label{tr:decorated}
\end{align}
For longer decorated words, define
\[
 Z_c[f_1,\ldots,f_k]
   =\sum_{n_1>\cdots>n_k\geq0}\prod_{j=1}^k f_j(n_j+c),
\]
and define $\HH_{a,b}[f_1,\ldots,f_k]$ by the same
prefix--suffix recursion as \eqref{scope:relative}; the displayed
formula \eqref{tr:decorated} is its depth-two case.
All decorated expressions below mean these initial definitions and their
specified continuations. They do not mean arbitrary interpolations of
finite sums. Write $f_s(x)=x^{-s}$ and use pointwise products of letters.

\begin{theorem}[The moving-index gap and every transverse derivative]
\label{tr:gap}
The two functions
\[
 A(s,u,t)=\HH_{a,b}[f_sQ(u,\cdot),f_t],\qquad
 B(s,u,t)=\HH_{a,b}[f_s,f_tQ(u,\cdot+1)]
\]
continue individually to entire functions of $(s,u,t)$. They obey
\begin{align}
 A(s,u,t)
   &=\HH_{a,b}(u,s,t)+\HH_{a,b}(s+u,t)
                      +Q(u,a)\HH_{a,b}(s,t),\label{tr:A}\\
 B(s,u,t)&=A(s,u,t)+\HH_{a,b}(s,u,t).\label{tr:B}
\end{align}
Thus
\begin{equation}\label{tr:genericgap}
 \boxed{\HH_{a,b}(s,u,t)
   =\HH_{a,b}[f_s,f_tQ(u,\cdot+1)]
       -\HH_{a,b}[f_sQ(u,\cdot),f_t].}
\end{equation}
For $m,r\geq0$ this gives
\begin{align}
 J_{m,r}(s,t;a,b)
 &:=\left.\partial_u^r\HH_{a,b}(s,u,t)\right|_{u=-m}\notag\\
 &=\HH_{a,b}[f_s,f_t\ell_{m,r}(\cdot+1)]
      -\HH_{a,b}[f_s\ell_{m,r}(\cdot),f_t].
 \label{tr:J}
\end{align}
In particular this is an identity for arbitrary independent $s,t$ and
arbitrary right-half-plane endpoints.
\end{theorem}

\begin{proof}
First replace $Q$ by $\zeta$ and take all real parts large.
In the decorated sum $Z_c[f_s\zeta(u,\cdot),f_t]$, the additional
tail index is either larger than the outer index or equal to it.
Accordingly
\[
 Z_c[f_s\zeta(u,\cdot),f_t]
          =Z_c(u,s,t)+Z_c(s+u,t),
 \qquad
 Z_c[f_s\zeta(u,\cdot)]=Z_c(u,s)+\zeta(s+u,c).
\]
Substitute these identities in \eqref{tr:decorated}. Expanding the
relative recursion \eqref{scope:relative} and collecting the same cuts
gives
\[
 \HH_{a,b}[f_s\zeta(u,\cdot),f_t]
   =\HH_{a,b}(u,s,t)+\HH_{a,b}(s+u,t)
                            +\zeta(u,a)\HH_{a,b}(s,t).
\]
In the other decorated double, its two indices above the lower one
can be ordered in either direction or can be equal. These three cases
give $Z_c(s,u,t)$, $Z_c(u,s,t)$, and $Z_c(s+u,t)$.
The same relative recursion gives the extra term $\HH_{a,b}(s,u,t)$.

Subtracting the constant letter $1/(u-1)$ proves
\eqref{tr:A}--\eqref{tr:B}. Their right sides are entire, so they also
provide individual continuations of the decorated doubles through every
spectral intersection. Their difference proves \eqref{tr:genericgap}.
Differentiate this entire identity. At $u=-m$ the constants that
distinguish the derivatives of $Q$ from $\ell_{m,r}$ cancel between
the two doubles. This proves \eqref{tr:J}.
\end{proof}

The arbitrary-depth version of the counting argument will be useful
when several unmarked slots have first been removed. It also makes the
general-endpoint meaning of every boundary gap explicit.

\begin{lemma}[A decorated word with one moving tail]
\label{tr:masterlemma}
Let $\boldsymbol s=(s_1,\ldots,s_k)$ and $1\leq i\leq k$.
Then the decorated relative word has the jointly entire continuation
\begin{align}
 &\HH_{a,b}[f_{s_1},\ldots,f_{s_i}Q(u,\cdot),\ldots,f_{s_k}]
 \notag\\
 &\quad=Q(u,a)\HH_{a,b}(s_1,\ldots,s_k)
   +\sum_{j=0}^{i-1}\HH_{a,b}(s_1,\ldots,s_j,u,s_{j+1},\ldots,s_k)
 \notag\\
 &\qquad\quad
   +\sum_{j=1}^{i}\HH_{a,b}(s_1,\ldots,s_j+u,\ldots,s_k).
 \label{tr:masterdecorate}
\end{align}
In particular the top and bottom boundary identities are
\begin{align}
 \HH_{a,b}(u,\boldsymbol s)
  &=\HH_{a,b}[f_{s_1}Q(u,\cdot+1),f_{s_2},\ldots,f_{s_k}]
                      -Q(u,a)\HH_{a,b}(\boldsymbol s),
 \label{tr:topgap}\\
 \HH_{a,b}(\boldsymbol s,u)
  &=Q(u,b)\HH_{a,b}(\boldsymbol s)
      -\HH_{a,b}[f_{s_1},\ldots,f_{s_k}Q(u,\cdot)].
 \label{tr:bottomgap}
\end{align}
\end{lemma}

\begin{proof}
First replace $Q(u,x)$ by $\zeta(u,x)$ and take all real parts
large. In the infinite ordered sum, the extra tail index is at least
the index in position $i$. Its possible strict positions are the $i$
slots preceding that position, and its possible equalities are the
first $i$ existing indices. Thus the infinite decorated sum is exactly
the sum of the insertion terms and merger terms in
\eqref{tr:masterdecorate}, with $Z_c$ in place of $\HH$ and
without the endpoint term.

Expand each such $Z_b$ by deconcatenation and group terms by the cut
in the original undecorated word. For a cut before position $i$,
insertions and mergers lying wholly in the prefix give
$\zeta(u,a)Z_a(\text{prefix})$ by the ordinary singleton stuffle
identity. The remaining insertions and mergers lie in the suffix.
For a cut at or after position $i$, the same counting identity gives
$Z_a(\text{decorated prefix})$. These are exactly the two types
of cuts in the defining recursion for the decorated relative word.
Triangular uniqueness proves the claimed identity with $\zeta$;
subtracting $\HH(\boldsymbol s)/(u-1)$ replaces its endpoint value
by $Q(u,a)$ and proves \eqref{tr:masterdecorate}. Its right side
provides the jointly entire continuation.

At $i=1$, use $Q(u,x+1)=Q(u,x)-x^{-u}$ to obtain
\eqref{tr:topgap}. At $i=k$, the relative singleton stuffle identity
combines all insertions and mergers, and
$Q(u,b)-Q(u,a)=h_{a,b}(u)$ gives \eqref{tr:bottomgap}.
The difference of the decorations at adjacent positions $i+1$ and $i$,
using $Q(u,\cdot+1)$ at the lower position, leaves exactly the
insertion at the gap between them. This proves the interior gap
identity at any depth as well. No integer-endpoint uniqueness principle
is used.
\end{proof}

At an integer interval $a=b+N$, the meaning of the gap is especially
transparent:
\begin{equation}\label{tr:finitegap}
 J_{m,r}(s,t;b+N,b)=
 \sum_{N>i>k\geq0}
 \frac{\ell_{m,r}(k+b+1)-\ell_{m,r}(i+b)}
      {(i+b)^s(k+b)^t}.
\end{equation}
Indeed the middle-index sum is
\[
 \sum_{k<j<i}(j+b)^m(-\log(j+b))^r
           =\ell_{m,r}(k+b+1)-\ell_{m,r}(i+b).
\]
The proof above establishes the general-endpoint identity independently
of this finite illustration.

The Hurwitz shift identity also gives an equivalent form using only
unshifted decorated letters:
\begin{align}
 J_{m,r}(s,t;a,b)
   ={}&\HH_{a,b}[f_s,f_t\ell_{m,r}]
       -\HH_{a,b}[f_s\ell_{m,r},f_t]
       -\partial_t^r\HH_{a,b}(s,t-m).
 \label{tr:commutator}
\end{align}
The last term is essential. For $m=r=0$, it recovers exactly
\[
 \HH(s,0,t)=\HH(s-1,t)-\HH(s,t-1)-\HH(s,t),
\]
the strict zero-gap deletion rule.

\subsection{Gamma gaps, Stieltjes primitives, and repeated primitives}

For the first transverse derivative at zero, \eqref{scope:hurwitz}
turns \eqref{tr:finitegap} into
\begin{equation}\label{tr:logGammafinite}
 \left.\partial_u\HH_{b+N,b}(s,u,t)\right|_{u=0}
 =
 \sum_{N>i>k\geq0}
 \frac{\log\Gamma(k+b+1)-\log\Gamma(i+b)}
      {(i+b)^s(k+b)^t}.
\end{equation}
The constants $\log(2\pi)/2$ disappear. This is the requested transverse
identity in explicit Gamma coordinates, with \eqref{tr:J} specifying
its general-endpoint counterpart.

Differentiating $-\!u\zeta(1+u,x)$ at $u=0$ gives, for $r\geq1$,
\begin{equation}\label{tr:Stieltjesprimitive}
 \frac{d}{dx}\ell_{0,r}(x)=(-1)^r r\,\gamma_{r-1}(x).
\end{equation}
Therefore
\begin{equation}\label{tr:integratedgap}
 \ell_{0,r}(z+1)-\ell_{0,r}(y)
   =(-1)^{r+1}r\int_{z+1}^{y}\gamma_{r-1}(x)\dd x.
\end{equation}
The path lies in the right half-plane. The integral is path independent,
and its additive normalization is fixed by the displayed endpoints.
Combining \eqref{tr:J} with \eqref{tr:integratedgap} gives all the
corresponding harmonic identities with Stieltjes-primitive gaps.
At $u=1$, the same entire identity instead gives
\begin{align}
 \left.\partial_u^r\HH(s,u,t)\right|_{u=1}
 =(-1)^r\bigl\{\HH[f_s,f_t\gamma_r(\cdot+1)]
                 -\HH[f_s\gamma_r,f_t]\bigr\}.
 \label{tr:Stieltjesgap}
\end{align}
These statements retain the ordered information; they do not claim
a reduction of every resulting value to depth-one constants.

A standard balanced normalization packages repeated primitives of
log-Gamma \cite{EspinosaMoll}:
\begin{equation}\label{tr:balanced}
 \mathcal B_m(x)=
       \frac{\zeta'(-m,x)+H_m\zeta(-m,x)}{m!},\qquad H_0=0.
\end{equation}
Here $\mathcal B_0(x)=\log\Gamma(x)-\tfrac12\log(2\pi)$, and
\begin{equation}\label{tr:balancedderivative}
 \mathcal B_m'(x)=\mathcal B_{m-1}(x)\qquad(m\geq1).
\end{equation}
To check this, differentiate the Hurwitz identity
$\partial_x\zeta(u,x)=-u\zeta(u+1,x)$ once in $u$, then use
$H_m=H_{m-1}+1/m$. The factor $m!$ in \eqref{tr:balanced} is
exactly the one needed for \eqref{tr:balancedderivative}.

It follows that the repeated-Gamma-primitive gap has the exact identity
\begin{align}
 &\HH[f_s,f_t\mathcal B_m(\cdot+1)]
                 -\HH[f_s\mathcal B_m,f_t]\notag\\
 &\qquad=\frac{J_{m,1}(s,t)+H_m\HH(s,-m,t)}{m!}.
 \label{tr:balancedgap}
\end{align}
The second term is a finite Bernoulli reduction. Thus no undetermined
polynomial or integration constant is hidden in the primitive ladder.

As one elementary Stieltjes check, let $h=h_{a,b}$. Differentiating the
stuffle of a singleton with $\HH(s,t)$ gives
\begin{align}
 &\left.\partial_u\{\HH(u,s,t)+\HH(s,u,t)+\HH(s,t,u)\}\right|_{u=0}
 \notag\\
 &\qquad=h'(0)\HH(s,t)-(\partial_s+\partial_t)\HH(s,t).
 \label{tr:insertionderivative}
\end{align}
At $s=t=1$, use $\HH(v,v)=(h(v)^2-h(2v))/2$ to obtain the fully
specified expression
\begin{equation}\label{tr:Stieltjesexample}
 \frac{h'(0)}2\{h(1)^2-h(2)\}
   +h(1)\{\gamma_1(b)-\gamma_1(a)\}
   +\zeta'(2,b)-\zeta'(2,a).
\end{equation}
This is a consistency specialization of the classical stuffle algebra,
not a separate claim of a new product principle.

\subsection{Absolutely convergent transverse Mellin formulas}

The following elementary normalized-Mellin lemma makes the transverse
construction directly computable.

\begin{lemma}[Every derivative at a nonpositive Mellin order]
\label{tr:Mellinlemma}
Suppose $f$ is smooth on $[0,\infty)$, extends holomorphically near zero,
and all its derivatives decrease exponentially at infinity. Let
\[
 \mathcal M_f(u)=\frac1{\Gamma(u)}
                  \int_0^\infty y^{u-1}f(y)\dd y,\qquad\Re u>0.
\]
It has an entire continuation. Put $c_j=[v^j]\Gamma(v)^{-1}$.
For $m\geq0$,
\begin{align}
 \mathcal M_f(-m)&=(-1)^m f^{(m)}(0),\label{tr:Mvalue}\\
 \mathcal M_f^{(r)}(-m)
 &=(-1)^m r!\sum_{j=1}^r\frac{c_j}{(r-j)!}
      \int_0^\infty\frac{\log^{r-j}y}{y}
          \{f^{(m)}(y)-f^{(m)}(0)e^{-y}\}\dd y ,
 \qquad r\geq1.
 \label{tr:Mjet}
\end{align}
All integrals in \eqref{tr:Mjet} are ordinary absolutely convergent
integrals.
\end{lemma}

\begin{proof}
Repeated integration by parts gives
$\mathcal M_f(u)=(-1)^m\mathcal M_{f^{(m)}}(u+m)$, initially on
$\Re u>0$ and then by continuation. It is enough to expand at zero.
Write $g=f^{(m)}$ and use the exact identity
\[
 \mathcal M_g(v)=g(0)+\frac1{\Gamma(v)}
       \int_0^\infty y^{v-1}\{g(y)-g(0)e^{-y}\}\dd y .
\]
The last integral is holomorphic for $\Re v>-1$, locally at zero.
Its $q$th derivative at zero is the integral with
$\log^qy/y$. The bracket is $O(y)$ near zero and decreases
exponentially at infinity. Coefficient multiplication by
$\Gamma(v)^{-1}$ proves \eqref{tr:Mjet}.
The continuation and value formula also follow from repeated integration
by parts with an arbitrarily large number of derivatives.
\end{proof}

The coefficients $c_j$ are explicit Gamma Taylor coefficients:
\[
 c_1=1,\qquad c_2=\gamma,\qquad
 c_3=\frac{\gamma^2-\zeta(2)}2,
\]
and in every degree they follow from
\[
 \frac1{\Gamma(v)}
     =v\exp\left(\gamma v+
             \sum_{q\geq2}\frac{(-1)^{q+1}\zeta(q)}qv^q\right).
\]

\begin{theorem}[Convergent formula for every transverse gap]
\label{tr:integral}
Let $\Re s,\Re t>0$, $m\geq0$, and $r\geq1$. In
\eqref{tr:Mjet}, set
\[
 f(y)=\mathscr K_{a,b}^{(3)}(x,y,z).
\]
Then $J_{m,r}(s,t;a,b)$ is the normalized $(x,z)$ Mellin transform
of its right side. Explicitly it is
\begin{align}
 &\frac{(-1)^m r!}{\Gamma(s)\Gamma(t)}
 \sum_{j=1}^r\frac{c_j}{(r-j)!}
 \int_0^\infty\int_0^\infty
       x^{s-1}z^{t-1}\notag\\
 &\qquad\cdot\int_0^\infty \frac{\log^{r-j}y}{y}
 \left\{\partial_y^m\mathscr K^{(3)}_{a,b}(x,y,z)
       -\partial_y^m\mathscr K^{(3)}_{a,b}(x,0,z)e^{-y}\right\}
             \dd y\dd z\dd x .
 \label{tr:tripleintegral}
\end{align}
All these integrals are absolutely convergent. Their common continuation
is entire in $s,t$. Every further fixed spectral derivative can be
taken in the convergent integrals on compact subsets of
$\Re s,\Re t>0$.
\end{theorem}

\begin{proof}
Lemma~\ref{zbg:bounds} gives uniform exponential control of all derivatives
of the kernel, including at $y=0$. The inner bracket is uniformly
$O(y)e^{-\delta(x+z)}$ as $y\downarrow0$; for $y\geq1$ it is bounded
by $Ce^{-\delta(x+z)}(e^{-\delta y}+e^{-y})$.
Consequently the inner logarithmic integrals and the two outer Mellin
integrals are absolutely convergent. Fubini's theorem and the preceding
lemma give the derivative of
$\mathscr E_{a,b}(s,u,t;\boldsymbol0)=\HH_{a,b}(s,u,t)$.
Theorem~\ref{zbg:entire} gives the entire continuation. Powers of $\log x$
and $\log z$ from further spectral derivatives preserve the same
integrable majorants on compact half-plane subsets.
\end{proof}

For $m=0,r=1$ the formula contains just
\[
 \int_0^\infty
 \frac{\mathscr K^{(3)}(x,y,z)-\mathscr K^{(3)}(x,0,z)e^{-y}}y\dd y.
\]
This supplies a convergent realization of the log-Gamma gap at arbitrary
noninteger endpoints, without separate divergent ordered sums.

\subsection{All negative outer values and fixed Barnes constants}

The weighted difference in Theorem~\ref{tr:jointD} has especially
explicit values at nonpositive outer orders.

\begin{theorem}[Finite Bernoulli formula for the transverse weighted tails]
\label{tr:negativeD}
Let $d,m,r\geq0$ and $q=d+1$. Then
\begin{align}
 D_{m,r}(-d;a,b)
   =\frac1q\biggl\{&
        B_q(a)\ell_{m,r}(a)-B_q(b)\ell_{m,r}(b)\notag\\
       &+\sum_{k=0}^q\binom qk B_{q-k}(1)
                             h_{a,b}^{(r)}(-m-k)\biggr\}.
 \label{tr:negativeformula}
\end{align}
Every term is an endpoint Hurwitz derivative or a Bernoulli polynomial.
\end{theorem}

\begin{proof}
The required top-index deletion identity is
\begin{equation}\label{tr:topdelete}
 \HH(-d,u)=\frac1q\left\{
 B_q(a)h(u)-\sum_{k=0}^q\binom qk B_{q-k}(1)h(u-k)\right\}.
\end{equation}
For a direct analytic derivation, sum the leading index of $Z_c(v,u)$
as $\zeta(v,n+c+1)$, then continue $v$ to $-d$ while $\Re u$ is
large. Use
$\zeta(-d,x+1)=-B_q(x+1)/q$ and expand the Bernoulli polynomial.
Insert the result in the relative quotient. The leading singleton
$Z_a(-d)=-B_q(a)/q$ contributes the first term. This proves
\eqref{tr:topdelete} on a nonempty domain, and entireness in $u$ extends it.

Differentiate \eqref{tr:topdelete} $r$ times in $u$ and set $u=-m$.
Substitute in \eqref{tr:Didentity} and use
$h(-d)=(B_q(a)-B_q(b))/q$. The terms involving
$B_q(a)\ell_{m,r}(b)$ cancel, giving \eqref{tr:negativeformula}.
\end{proof}

For $m=0,r=1,d=0$, the standard Barnes relation
\cite{Vardi,DLMFBarnes} gives the particularly simple identity
\begin{equation}\label{tr:Barnes}
 \boxed{D_{0,1}(0;a,b)=
    \log G(a)-\log G(b)-\frac{a-b}{2}\log(2\pi).}
\end{equation}
Here $G$ is Barnes' function with its standard asymptotic normalization,
$G(1)=1$, and $G(x+1)=\Gamma(x)G(x)$. Holomorphic logarithms in the
right half-plane are normalized from the positive real axis.
Indeed the $q=1$ formula is
\[
 (a-1)\ell_{0,1}(a)-(b-1)\ell_{0,1}(b)
                    +\zeta'(-1,b)-\zeta'(-1,a),
\]
and
$\log G(x)=(x-1)\log\Gamma(x)+\zeta'(-1)-\zeta'(-1,x)$
proves \eqref{tr:Barnes}. At an integer endpoint difference it also
agrees with the finite product obtained by iterating the functional
equation. General $d$ gives the analogous finite endpoint formula in
higher Hurwitz derivatives, without choosing an unspecified multiple
Gamma normalization.

\subsection{A finite family after arbitrary unmarked deletions}

The preceding construction is stable under inserting any finite word
of unmarked nonpositive indices. We state exactly the scope.

\begin{proposition}[One marked index with arbitrary nonpositive decoration]
\label{tr:finitefamily}
Suppose a word contains $k\geq1$ independent free indices, one marked
index $u$, and unmarked indices $-m_1,\ldots,-m_e$, with $m_i\geq0$.
Put $M=\sum_i(m_i+1)$. After differentiating the marked index $r$ times
and setting $u=-m$, the word is a finite polynomial-endpoint combination
of $k$-letter relative sums whose letters are ordinary shifted powers,
ordinary spectral derivatives of those powers, and at most one factor
\[
 \ell_{m+\nu,r}(x),\qquad 0\leq\nu\leq M.
\]
Endpoint values of the same finite family may occur. Thus the new gap
functions needed for this one-marked problem are a specified finite
list. This is a bound for this representation, not an assertion of
minimal arithmetic depth.
\end{proposition}

\begin{proof}
First regard $u$ as free. The finite Bernoulli deletion rules for each
unmarked $-m_i$ hold identically in all remaining indices: sum its
finite strict gap by a Bernoulli polynomial, or use the leading-tail
argument of \eqref{tr:topdelete} at the upper boundary, and insert the
result in the relative recursion. An interior deletion has the two
endpoint terms
\[
 \frac1{m_i+1}\sum_{j=0}^{m_i+1}\binom{m_i+1}{j}
 \left\{B_{m_i+1-j}\HH(\ldots,s-j,t,\ldots)
       -B_{m_i+1-j}(1)\HH(\ldots,s,t-j,\ldots)\right\}.
\]
It shifts adjacent surviving orders down by nonnegative integers.
The total shift budget is at most $M$, even when adjacent unmarked
indices are combined during the deletion. Coefficients are endpoint
polynomials and are independent of every surviving free order.

Now differentiate in $u$ and set $u=-m$. Each surviving marked order
is $-m-\nu$ for some $0\leq\nu\leq M$.
Apply Lemma~\ref{tr:masterlemma} at the surviving marked position.
Its adjacent-decoration difference gives the interior gap, and
\eqref{tr:topgap}--\eqref{tr:bottomgap} give the boundary gaps.
Their finite-interval forms are
$\ell(y+1)-\ell(a)$ at the top and
$\ell(b)-\ell(y)$ at the bottom. Only one decorated factor is introduced.
Finally
$\ell_{m+\nu,r}(x+1)=\ell_{m+\nu,r}(x)
                   -x^{m+\nu}(-\log x)^r$
removes shifted potentials in favor of ordinary free-order derivatives.
The analytic identities at every step were proved with the marked
order still free, so the transverse differentiation is legitimate.
\end{proof}

There is also an analytic all-decoration version. Take any fixed
free word with one marked slot, apply Theorem~\ref{zbg:main}, then
differentiate that slot at $-m$. Local normality proves convergence on
the same full unit polydisc. Formula \eqref{tr:Mjet}, applied to the
marked kernel coordinate, gives the corresponding logarithmic-integral
generator. This simultaneously extends R7 and R8 to every derivative
order and every nonpositive marked center.

If there are no other free indices, the remaining singleton is
$h_{a,b}^{(r)}(-m-\nu)$, so no decorated lower-depth sum is needed.
Several independently marked indices are a further problem: this
proposition makes no claim that their repeated gap primitives close
in the same one-marked list.


% End sections/03_transverse.tex

% Begin sections/04_mixed_moments.tex
\section{Mixed centered harmonic-tail moments}
\label{mt:section}

The source reports evaluate every centered even moment of a trigamma
square \cite{IncomingHigherReflection,IncomingParity}; the harmonic-parity
report asks for the corresponding evaluation for two unequal even-order
harmonic tails. We resolve that two-factor problem for all
orders and all moments. The resulting finite formula reduces entirely to
ordinary zeta values and their products. We then package the infinite
moment sequence into one analytic generator with explicit
polylogarithmic primitives. The centered parity mechanism is inherited
from the source criterion; the evaluation and the generator are the
additional results here. Euler's classical odd-weight reduction supplies
the arithmetic step~\cite{Bradley2007}.

\subsection{Continuation and ordinary convergent representatives}

Let $a,b\ge2$ be even integers, and put
\begin{equation}\label{mt:notation}
 T_a(n)=\zeta(a,n+1)=\zeta(a)-H_n^{(a)}
       =\frac{\psi^{(a-1)}(n+1)}{(a-1)!},
 \qquad x_n=n+\frac12.
\end{equation}
The last equality uses that $a$ is even. Define
\begin{equation}\label{mt:F}
 F_{a,b}(s)=\sum_{n=0}^{\infty}\frac{T_a(n)T_b(n)}{x_n^s},
 \qquad V_{a,b}(M)=F_{a,b}(-2M),\quad M\ge0.
\end{equation}
The series for $F_{a,b}$ converges absolutely when
$\Re s>3-a-b$. The value at $-2M$ is understood through the continuation
proved below whenever the unmodified series diverges.

Write $\alpha=a+b-2$, and define rational coefficients
\begin{align}
 t_{a,0}&=\frac1{a-1},&
 t_{a,k}&=\frac{B_{2k}(1/2)}{(2k)!}(a)_{2k-1}\quad(k\ge1),
 &c_k&=\sum_{j=0}^k t_{a,j}t_{b,k-j}.
 \label{mt:centered-coefficients}
\end{align}
Here $(a)_r$ is the rising factorial, $B_j(x)$ is the Bernoulli
polynomial, and $B_j=B_j(0)$, with $B_1=-1/2$.
The centered Hurwitz expansions \cite[\S25.11(xii)]{DLMFHurwitz} give
\begin{equation}\label{mt:centered-expansion}
 T_a(n)\sim\sum_{k\ge0}t_{a,k}x_n^{1-a-2k},
 \qquad
 T_a(n)T_b(n)\sim\sum_{k\ge0}c_kx_n^{-\alpha-2k}.
\end{equation}
Only finite truncations of these asymptotic expansions are used.

\begin{proposition}[Regularity and bracketed moments]
\label{mt:regularity}
The function $F_{a,b}$ extends meromorphically to the complex plane.
Its only possible poles are simple poles at
$s=3-a-b-2k$, $k\ge0$, with residue $c_k$.
In particular every even spectral integer is regular. For
\begin{equation}\label{mt:subtraction-polynomial}
 P_M(x)=\sum_{\alpha+2k\le2M}c_kx^{2M-\alpha-2k},
\end{equation}
one has the ordinary convergent identity
\begin{equation}\label{mt:bracket}
 V_{a,b}(M)=\sum_{n=0}^{\infty}
 \left[x_n^{2M}T_a(n)T_b(n)-P_M(x_n)\right].
\end{equation}
The polynomial is zero if its indexing set is empty. The bracket is
$O(x_n^{-2})$.
\end{proposition}

\begin{proof}
For any fixed $K$, subtract the first $K+1$ terms of
\eqref{mt:centered-expansion} inside the defining series. The remainder
converges normally in a right half-plane extending farther to the left,
while the removed terms
sum to
$\sum_{k=0}^K c_k\zeta(s+\alpha+2k,1/2)$.
Increasing $K$ provides compatible continuations and the stated poles
and residues. Since $a+b$ is even, all these potential poles are odd
integers. At $s=-2M$, every nondecaying term has a nonnegative even
power of $x_n$. Its continued sum is
$\zeta(-2j,1/2)=0$, including $j=0$. After their removal the next
possible power is $x_n^{-2}$, which proves \eqref{mt:bracket}.
\end{proof}

\begin{corollary}[Every spectral derivative with its correction terms]
\label{mt:spectralderivatives}
For $r,M\geq0$, every derivative at a centered even spectral value has
the ordinary convergent representation
\begin{align}
 F_{a,b}^{(r)}(-2M)
  ={}&\sum_{n=0}^{\infty}(-\log x_n)^r
       \left[x_n^{2M}T_a(n)T_b(n)-P_M(x_n)\right]\notag\\
    &+\sum_{\alpha+2k\leq2M}
               c_k\zeta^{(r)}(\alpha+2k-2M,1/2).
 \label{mt:logmoment}
\end{align}
The second sum vanishes for $r=0$. It must in general be retained for
$r\geq1$.
\end{corollary}
\begin{proof}
Use the same finite subtractions as in \eqref{mt:subtraction-polynomial}
at a moving $s$ in a neighborhood of $-2M$. The remainder series and all
its fixed spectral derivatives converge locally normally there, since
their terms are bounded by $C x_n^{-2+\varepsilon}(1+|\log x_n|^r)$
with some $\varepsilon<1$. Differentiation of the finite Hurwitz
correction gives exactly \eqref{mt:logmoment}.
\end{proof}

For example, the first derivative in the $(2,4)$ family at $s=-4$ is
\begin{equation}\label{mt:derivative-example}
 F'_{2,4}(-4)=
 -\sum_{n\geq0}\log x_n
        \left[x_n^4T_2(n)T_4(n)-\frac13\right]
 -\frac{\log 2}{6},
\end{equation}
because $\zeta'(0,1/2)=-\tfrac12\log2$. This identity fixes the
logarithmic remainder and its constant; it does not claim that the
remainder is an ordinary-zeta combination.

We will also use a concrete equivalent characterization. Let
$\operatorname{CT}_N$ mean the coefficient of $N^0(\log N)^0$ in
an asymptotic expansion in the integer cutoff $N$. Then
\begin{equation}\label{mt:cutoff-value}
 V_{a,b}(M)=\operatorname{CT}_N
   \sum_{n=0}^{N-1}x_n^{2M}T_a(n)T_b(n).
\end{equation}
Indeed the partial sum of each removed even polynomial power is
$B_{2j+1}(N+1/2)/(2j+1)$, an odd polynomial in $N$ with zero constant
term. The summable remainder in \eqref{mt:bracket} tends to its ordinary
sum. The coordinate $N$ in \eqref{mt:cutoff-value} is part of the
statement; no change of endpoint coordinate is implicit.

\subsection{A finite formula at every moment order}

The double-zeta convention used below is
\[
 \zeta(u,v)=\sum_{j>k\ge1}\frac1{j^u k^v}.
\]
Thus the first exponent belongs to the larger summation variable.
For each positive odd integer $d$, define
\begin{equation}\label{mt:oriented}
 K_{a,b}(d)=
 \begin{cases}
  \displaystyle\zeta(b,a-d)+\frac12\zeta(a+b-d),&d<a,\\[2mm]
  \displaystyle\sum_{j=1}^{\min(b/2,(d-a+1)/2)}
    \frac{\binom{d-a+1}{2j}B_{d-a+1-2j}}{d-a+1}
             \zeta(b-2j),&d>a,
 \end{cases}
\end{equation}
and put
\begin{equation}\label{mt:E}
 E_{a,b}(d)=K_{a,b}(d)+K_{b,a}(d).
\end{equation}
There is no missing case $d=a$: $d$ is odd and $a$ even. All double
zetas in the first branch are ordinary convergent sums. In the second
branch, $\zeta(0)=-1/2$ has its usual continued meaning.

\begin{theorem}[All mixed even-tail moments]
\label{mt:all-moments}
For even $a,b\ge2$ and every integer $M\ge0$,
\begin{equation}\label{mt:moment-formula}
 \boxed{\displaystyle
 V_{a,b}(M)=\frac1{2M+1}\sum_{\ell=0}^{M}
  \binom{2M+1}{2\ell+1}B_{2M-2\ell}(1/2)
                          E_{a,b}(2\ell+1).}
\end{equation}
The summation and every coefficient in this formula are finite and
explicit.
\end{theorem}

The $j=b/2$ term in \eqref{mt:oriented}, whenever its range includes
it, records an exact rational boundary contribution through
$\zeta(0)=-1/2$. Discarding this term would already give an incorrect
answer for divergent centered moments of the trigamma square.

\begin{proof}
First form the exact polynomial partial sum
\begin{align}
 S_M(k)&=\sum_{n=0}^{k-1}(n+1/2)^{2M}
       =\frac{B_{2M+1}(k+1/2)}{2M+1}
       =\sum_{\ell=0}^{M}s_{M,\ell}k^{2\ell+1},
 \label{mt:S}\\
 s_{M,\ell}&=\frac{\binom{2M+1}{2\ell+1}}{2M+1}
                         B_{2M-2\ell}(1/2).
 \notag
\end{align}
The tail recurrence gives
\begin{align*}
&T_a(k-1)T_b(k-1)-T_a(k)T_b(k)\\
&\hspace{18mm}=k^{-a}T_b(k-1)+k^{-b}T_a(k-1)-k^{-a-b}.
\end{align*}
Applying finite summation by parts, without taking any divergent limits,
therefore yields
\begin{equation}\label{mt:finite-sbp}
 \sum_{n=0}^{N-1}x_n^{2M}T_a(n)T_b(n)
 =\sum_{\ell=0}^{M}s_{M,\ell}\,\mathcal E_{2\ell+1}(N),
\end{equation}
where
\begin{align}
 \mathcal E_d(N)={}&N^dT_a(N)T_b(N)
   +\sum_{k=1}^N k^{d-a}T_b(k-1)
   +\sum_{k=1}^N k^{d-b}T_a(k-1)
   -H_N^{(a+b-d)}.
 \label{mt:finite-E}
\end{align}
We compute its cutoff constant explicitly.

If $d<a$, the first oriented sum in \eqref{mt:finite-E} is absolutely
convergent, and direct rearrangement gives
\begin{equation}\label{mt:convergent-oriented}
 \sum_{k=1}^{\infty}k^{d-a}T_b(k-1)
 =\zeta(b,a-d)+\zeta(a+b-d).
\end{equation}
If $d>a$, put $p=d-a$, a positive odd integer. The Faulhaber
polynomial is
\begin{equation}\label{mt:faulhaber}
 Q_p(X)=\sum_{k=1}^X k^p
   =\frac12X^p+\sum_{j=1}^{(p+1)/2}q_{p,2j}X^{2j},
 \qquad
 q_{p,2j}=\frac{\binom{p+1}{2j}B_{p+1-2j}}{p+1}.
\end{equation}
The finite-sum notation specifies the polynomial by its values at
nonnegative integers. Define $q_{p,h}=0$ outside its displayed
nonzero even degrees, whenever an even $h$ is used below. A finite
rearrangement which retains the infinite tail gives the exact equality
\begin{equation}\label{mt:finite-rearrangement}
 \sum_{k=1}^N k^pT_b(k-1)
  =\sum_{h\ge1}[X^h]Q_p(X)\,H_N^{(b-h)}
    +Q_p(N)T_b(N).
\end{equation}
For $v>1$, the cutoff constant of $H_N^{(v)}$ is $\zeta(v)$;
for $v=1$ it is Euler's constant $\gamma$; and for integer $v\le0$
it is zero because the finite power sum is a polynomial in $N$ with
zero constant term.

At the potential logarithmic collision $d=a+b-1$, the half-leading
monomial in each of the two Faulhaber polynomials contributes
$\gamma/2$. The diagonal term in \eqref{mt:finite-E} contributes
$-\gamma$. These cancel. In particular, this proof never evaluates
an isolated $\zeta(1)$.

We now evaluate the constant contributions of the explicit tails. Write
\begin{equation}\label{mt:tau}
 T_a(N)\sim\sum_j\tau_a(j)N^{-j},\qquad
 \begin{cases}
 \tau_a(a-1)=1/(a-1),\\
 \tau_a(a)=-1/2,\\
 \tau_a(a+2k-1)=B_{2k}(a)_{2k-1}/(2k)!,&k\ge1,
 \end{cases}
\end{equation}
with all other coefficients zero. Since $a,b$ are even and $d$ odd,
an odd total degree in the product can only use one exceptional even
index $a$ or $b$. Hence
\begin{equation}\label{mt:boundary-product}
 \operatorname{CT}_N N^dT_a(N)T_b(N)
 =-\frac12\tau_b(d-a)-\frac12\tau_a(d-b).
\end{equation}
For $p=d-a>0$, the only odd monomial of $Q_p$ is $X^p/2$,
and the only even index in $T_b$ is $b$. Therefore
\begin{equation}\label{mt:boundary-oriented}
 \operatorname{CT}_N Q_p(N)T_b(N)
 =\frac12\tau_b(p)-\frac12q_{p,b}.
\end{equation}
The $\tau$ terms in \eqref{mt:boundary-product} cancel those from
the two instances of \eqref{mt:boundary-oriented}. The remaining
constant is
\begin{equation}\label{mt:boundary-constant}
 -\frac12\bigl(q_{d-a,b}+q_{d-b,a}\bigr),
\end{equation}
where a term is zero if its first index is not positive. This is
exactly the contribution of the $\zeta(0)$ terms in
\eqref{mt:oriented}.

For completeness consider the remaining diagonal constant. If
$W=a+b-d\ge3$, a branch $d<a$ contributes the full $\zeta(W)$
in \eqref{mt:convergent-oriented}, and a branch $d>a$ contributes
$\zeta(W)/2$ from the half-leading term in
\eqref{mt:finite-rearrangement}. After subtracting the diagonal
$\zeta(W)$, exactly one half remains for each branch $d<a$.
When $W=1$ the constants cancel as already shown. When $W<1$,
the diagonal power sum has zero constant term. The other surviving
Faulhaber terms have even degree and yield the even zeta values in
\eqref{mt:oriented}. Thus
\[
 \operatorname{CT}_N\mathcal E_d(N)=E_{a,b}(d).
\]
Taking constants in \eqref{mt:finite-sbp} and applying
\eqref{mt:cutoff-value} proves \eqref{mt:moment-formula}.
\end{proof}

\subsection{Reduction to ordinary zeta values}

For clarity, the classical reduction needed to eliminate every double
coordinate is written out. Let $m\ge2$ be even, let $n\ge1$ be odd,
and put $w=m+n$ and
$\mathcal U_{m,n}=\zeta(m,n)+\zeta(w)/2$. Euler's formulas, in the
convention above, give
\begin{equation}\label{mt:euler-one}
 \mathcal U_{m,1}=\frac{m+1}{2}\zeta(m+1)
  -\sum_{j=1}^{m/2-1}\zeta(2j)\zeta(m+1-2j),
\end{equation}
and, for $n\ge3$,
\begin{align}
 \mathcal U_{m,n}={}&\frac12\binom wm\zeta(w)+\zeta(m)\zeta(n)
 \notag\\
 &-\sum_{j=1}^{(w-3)/2}
  \left[\binom{w-2j-1}{m-1}+\binom{w-2j-1}{n-1}\right]
                      \zeta(2j)\zeta(w-2j).
 \label{mt:euler-odd}
\end{align}
These are the ordinary specialization of the classical signed Euler
reduction; an elementary proof from partial fractions is provided in
\cite[Eq.~(2)--(3) and Proposition~1]{Bradley2007}. No divergent
$\zeta(1)$ appears in \eqref{mt:euler-one}--\eqref{mt:euler-odd}.
The first branch of \eqref{mt:oriented} is exactly
$\mathcal U_{b,a-d}$, so the formulas apply to every such branch.

\begin{corollary}[Ordinary-zeta closure and a fixed spanning list]
\label{mt:closure}
Every $V_{a,b}(M)$ is a rational linear combination of $1$, ordinary
even zeta values, ordinary odd zeta values, and products of one even
and one odd zeta value. Put $B=\max(a,b)$. The entire moment sequence
lies in the rational span of the at most $B$ coordinates
\begin{equation}\label{mt:finite-span}
 E_{a,b}(1),E_{a,b}(3),\ldots,E_{a,b}(B-1),
 \qquad 1,\zeta(2),\zeta(4),\ldots,\zeta(B-2).
\end{equation}
The positive even-zeta sublist is empty when $B=2$.
\end{corollary}
\begin{proof}
Ordinary-zeta closure follows by inserting
\eqref{mt:euler-one}--\eqref{mt:euler-odd} into
\eqref{mt:oriented}. For odd $d>B$, both branches in
\eqref{mt:oriented} are finite combinations of
$\zeta(0),\zeta(2),\ldots,\zeta(B-2)$. Formula
\eqref{mt:moment-formula} then proves the spanning assertion.
\end{proof}

The corollary is a finite spanning theorem, not a claim that its
coordinates are arithmetically independent or minimal. It answers the
ordinary-constant reduction question for every pair of even-order
tails; increasing the moment order introduces no additional constants.

\subsection{The first unequal-order family}

For $a=2$ and $b=4$, the first five coordinates are
\begin{align}
 E_{2,4}(1)&=\frac{15}{2}\zeta(5)-3\zeta(2)\zeta(3),
 &E_{2,4}(3)&=\frac12\zeta(2)+\frac32\zeta(3),\notag\\
 E_{2,4}(5)&=\frac14\zeta(2)-\frac38,
 &E_{2,4}(7)&=-\frac1{12}\zeta(2)-\frac13,\label{mt:E24}\\
 E_{2,4}(9)&=\frac1{12}\zeta(2)+\frac3{16}.\notag
\end{align}
The centered product expansion begins
\begin{equation}\label{mt:T24-expansion}
 T_2(n)T_4(n)=\frac1{3x_n^4}-\frac7{36x_n^6}
                       +\frac{61}{360x_n^8}+O(x_n^{-10}).
\end{equation}
Theorem~\ref{mt:all-moments} consequently gives the following ordinary
convergent identities:
\begin{align}
 \sum_{n\ge0}T_2(n)T_4(n)
  &=\frac{15}{2}\zeta(5)-3\zeta(2)\zeta(3),
 \label{mt:V0}\\
 \sum_{n\ge0}x_n^2T_2(n)T_4(n)
  &=\frac14\zeta(2)\zeta(3)+\frac16\zeta(2)+\frac12\zeta(3)
                   -\frac58\zeta(5),
 \label{mt:V1}\\
 \sum_{n\ge0}\left[x_n^4T_2(n)T_4(n)-\frac13\right]
  &=-\frac7{80}\zeta(2)\zeta(3)-\frac1{30}\zeta(2)
      -\frac14\zeta(3)+\frac7{32}\zeta(5)-\frac3{40}.
 \label{mt:V2}
\end{align}
One further example is
\begin{align}
 &\sum_{n\ge0}\left[x_n^6T_2(n)T_4(n)-\frac{x_n^2}{3}+\frac7{36}\right]
 \notag\\
 &\qquad=\frac{31}{448}\zeta(2)\zeta(3)-\frac1{672}\zeta(2)
          +\frac7{32}\zeta(3)-\frac{155}{896}\zeta(5)
          +\frac{31}{672}.
 \label{mt:V3}
\end{align}
Since $T_2=\psi'$ and $T_4=\psi'''/6$, multiplication by $6$
turns these into identities for products of two unequal polygamma
orders. They equally evaluate the centered harmonic numerator
$(H_n^{(2)}-\zeta(2))(H_n^{(4)}-\zeta(4))$.

\subsection{An analytic generator for all moment orders}

For a fixed ordered pair $(a,b)$, define
\begin{equation}\label{mt:G}
 G_b(t)=\frac12\coth(t/2)
             \sum_{j=1}^{b/2}\zeta(b-2j)\frac{t^{2j}}{(2j)!}.
\end{equation}
Its apparent singularity at zero is removable. It is odd and analytic
on $|t|<2\pi$.

\begin{theorem}[All-moment generating function]
\label{mt:generator}
Let $A_{a,b}$ be the unique analytic solution on $|t|<2\pi$ of
\begin{equation}\label{mt:generator-ode}
 A_{a,b}^{(a)}(t)=G_b(t),
\end{equation}
whose Taylor polynomial of degree less than $a$ is
\begin{equation}\label{mt:initial-data}
 \sum_{\substack{1\le d<a\\d\text{ odd}}}
  \left[\zeta(b,a-d)+\frac12\zeta(a+b-d)\right]\frac{t^d}{d!}.
\end{equation}
Then the exponential series
\begin{equation}\label{mt:moment-egf}
 \mathscr V_{a,b}(t)=\sum_{M=0}^{\infty}
                      V_{a,b}(M)\frac{t^{2M}}{(2M)!}
\end{equation}
converges for $|t|<2\pi$ and satisfies
\begin{equation}\label{mt:generator-formula}
 \boxed{\displaystyle
 \mathscr V_{a,b}(t)=
       \frac{A_{a,b}(t)+A_{b,a}(t)}{2\sinh(t/2)}.}
\end{equation}
The quotient has a removable value at zero and is even analytic.
\end{theorem}

\begin{proof}
Anchored integration of the analytic function $G_b$ gives existence
and uniqueness. Its oddness, and the even value of $a$, show that
$A_{a,b}$ is odd. The Bernoulli generating identity
\[
 C(t)=\frac t2\coth(t/2)
       =\sum_{k\ge0}\frac{B_{2k}}{(2k)!}t^{2k}
\]
gives
\[
 G_b(t)=C(t)\sum_{j=1}^{b/2}\zeta(b-2j)\frac{t^{2j-1}}{(2j)!}.
\]
For positive odd $p$, its coefficient of $t^p/p!$ is
\[
 \sum_{j=1}^{\min(b/2,(p+1)/2)}
   \frac{\binom{p+1}{2j}B_{p+1-2j}}{p+1}\zeta(b-2j).
\]
Together with the initial data, this says exactly that
\begin{equation}\label{mt:A-series}
 A_{a,b}(t)=\sum_{\substack{d\ge1\\d\text{ odd}}}
                             K_{a,b}(d)\frac{t^d}{d!}.
\end{equation}
Finally,
\[
 \frac1{2\sinh(t/2)}
     =\frac1t\sum_{k\ge0}\frac{B_{2k}(1/2)}{(2k)!}t^{2k}.
\]
Multiplication with the sum of the two series in
\eqref{mt:A-series} gives precisely the coefficient formula
\eqref{mt:moment-formula}. The numerator is odd and vanishes at zero;
the denominator has a simple zero there and no other zero in this
disk. Hence the quotient is even analytic throughout the disk, proving
actual convergence of \eqref{mt:moment-egf} as well as the identity.
\end{proof}

This analytic generator concerns the regularized values already fixed
by \eqref{mt:bracket}. It is not obtained by interchanging an exponential
series with the divergent unmodified moment sums.

\subsection{Explicit polylogarithmic primitives and their constants}

For an integer $k\ge1$, set
\begin{equation}\label{mt:L-integral}
 L_k(t)=\int_0^t\frac{u^k}{e^u-1}\,du.
\end{equation}
The integrand is analytic at zero. Initially in $\Re t>0$, the standard
polylogarithm differentiation rule \cite[\S25.12]{DLMFPolylog} gives the explicit anchored formula
\begin{equation}\label{mt:L-polylog}
 L_k(t)=k!\zeta(k+1)
 -\sum_{h=0}^k\frac{k!}{(k-h)!}
                      t^{k-h}\operatorname{Li}_{h+1}(e^{-t}).
\end{equation}
For real $t>0$, all polylogarithms have arguments in $(0,1)$; this fixes
the branch. The formula extends to the right half-plane by analyticity.
Differentiating the finite sum makes its adjacent terms cancel and
leaves $t^k/(e^t-1)$. As $t\to0$ from the right, all positive powers
multiplying the polylogarithms vanish, including
$t^k\operatorname{Li}_1(e^{-t})$, and the remaining term cancels
$k!\zeta(k+1)$. This proves both the derivative and the anchoring
constant. The completed expression, taken as a whole, then continues
analytically to $|t|<2\pi$ through its integral definition. Individual
polylogarithm terms need not be assigned separate analytic values
around the logarithmic point $t=0$.

For positive even $m$, define
\begin{align}
 R_{a,m}(t)={}&\frac{t^{a+m}}{2(a+m)!}
 \notag\\
 &+\frac1{m!(a-1)!}\sum_{r=0}^{a-1}(-1)^r\binom{a-1}{r}
                                  t^{a-1-r}L_{m+r}(t).
 \label{mt:R}
\end{align}
Expanding the anchored integration kernel $(t-u)^{a-1}$ shows that
\begin{equation}\label{mt:R-integral}
 R_{a,m}(t)=\frac1{(a-1)!}\int_0^t(t-u)^{a-1}
                         \frac12\coth(u/2)\frac{u^m}{m!}\,du.
\end{equation}
Here one uses $\tfrac12\coth(u/2)=\tfrac12+(e^u-1)^{-1}$;
the first term gives the polynomial in \eqref{mt:R}, and the second
gives the finite sum of $L$ integrals. Consequently
\begin{equation}\label{mt:A-polylog}
 A_{a,b}(t)=(\text{the polynomial in \eqref{mt:initial-data}})
              +\sum_{j=1}^{b/2}\zeta(b-2j)R_{a,2j}(t).
\end{equation}
Equations \eqref{mt:euler-one}--\eqref{mt:euler-odd} and
\eqref{mt:generator-formula}--\eqref{mt:A-polylog} express the whole
moment sequence using ordinary zeta constants, elementary functions,
and finitely many $\operatorname{Li}_k(e^{-t})$ with $k\le a+b$.
Every repeated primitive is anchored at zero. Near zero, the ODE or
Taylor form is more suitable for numerical evaluation because the
separate polylogarithmic terms undergo strong cancellation.

\subsection{Scope of the result}

The theorem resolves the two-factor even-tail sector of the source
question. It does not assert ordinary-zeta closure for the first
spectral derivative $F'_{a,b}(-2M)$, independent order derivatives,
or products of four tails. In particular, adding a spectral derivative
introduces a global logarithmic moment; the finite-value theorem alone
does not eliminate it. Verification records and precise next questions
are given in Sections~\ref{audit:section} and~\ref{research:section}.

% End sections/04_mixed_moments.tex

% Begin sections/05_audit_verification.tex
\Needspace{10\baselineskip}
\section{Source audit, normalization safeguards, and verification}
\label{audit:section}

\subsection{Inspected baseline and source attribution}

The source anchor is the ProveIt revision
\src{7bd45777a8a127e5dc69aff8887ca36a79a3059d}, including the eleven
archives then present in \src{docs/incoming} \cite{Repository}.
The archive inventory, immutable source URLs, and hashes accompany
the delivery.  The mathematical review concentrated on the relative
Hurwitz definitions, nonpositive-index reductions, one-free-order
generator, centered moment formulas, and the explicit R1, R7, and R8
questions of the incoming harmonic-parity report
\cite{IncomingParity,IncomingIndependent}.  This is a targeted analytic
review with a complete archive inventory; it is not a line-by-line
certification of every proof in the entire repository.

No newly confirmed false theorem in the inspected source material
emerged from this review.  The refinements below specify hypotheses and
normalizations that must survive integration.  In particular, the
source's warning against differentiating a fixed deleted index was
correct: Section~\ref{tr:section} supplies the extra moving-parameter
identity needed to cross that boundary.

The kernel attribution deserves explicit preservation.  The geometric
kernel reorganized in \eqref{zbg:kernel}, and the original entire
interpolation that it represents, are established in Ihara, Nakamura,
and Yamamoto \cite[Proposition~3.3 and Theorem~4.1]{INY}.  The published
reference is \emph{The Ramanujan Journal} \textbf{67} (2025), article~1,
doi:10.1007/s11139-025-01045-2; the arXiv version remains a useful
accessible precursor.  Bibliographic entries that describe this work
only as a preprint can therefore be updated.  The incoming
independent-orders report already identifies the interpolant
\cite{IncomingIndependent}.  Our contribution is the stated joint
twist completion, the all-block unit-polydisc theorem, and the
transverse formulas, rather than a new discovery of their underlying
geometric interpolation.

Likewise, the finite gap polynomials and the one-free-order generating
function retain their attribution to \cite{IncomingParity}.  Euler's
odd-weight double-zeta reduction remains a classical ingredient,
with the explicit primary proof cited in \cite{Bradley2007}.
Several incoming reports have overlapping parity and quartic-zeta
results from earlier common baselines.  Their separate proofs may be
useful, but the canonical question ledger should count the corresponding
resolved question once and retain each report's actual scope.

\subsection{An absolutely convergent boundary series can have a different value}
\label{audit:boundary}

The initial convergence domain in Theorem~\ref{tr:jointD} is essential.
An isolated exponent can cancel the leading asymptotic term of a
paired summand, producing an ordinary convergent series whose sum differs
from the continuation of the original family.  Here is a complete
elementary example, requiring no numerical regularization.

Take $m=r=0$, $a=2$, and $b=1$.  Then
\[
 \ell_{0,0}(x)=\zeta(0,x)=\frac12-x,
 \qquad h_{2,1}(s)=1.
\]
On the initial half-plane $\Re s>2$, the defining weighted difference
is
\begin{align}
 D_{0,0}(s;2,1)
 &=\sum_{n\geq0}\{f_s(n+1)-f_s(n+2)\},
 & f_s(x)&=\frac12x^{-s}-x^{1-s}.
 \label{audit:telescoping}
\end{align}
Its $N$th partial sum is exactly
\begin{equation}\label{audit:boundary-partial}
 -\frac12-\frac12(N+1)^{-s}+(N+1)^{1-s}.
\end{equation}
Consequently its entire continuation is the constant
\begin{equation}\label{audit:continued}
                         D_{0,0}(s;2,1)=-\frac12.
\end{equation}
But substituting $s=1$ in the individual paired summands first gives
the ordinary absolutely convergent series
\begin{equation}\label{audit:literal}
 \sum_{n\geq0}\frac12
       \left(\frac1{n+1}-\frac1{n+2}\right)=\frac12.
\end{equation}
The values differ by one, which is the boundary term
$(N+1)^{1-s}$ in \eqref{audit:boundary-partial}.  This discrepancy
does not contradict analytic continuation: the paired series lacks
locally normal convergence on a complex neighborhood of $s=1$.

Thus \eqref{tr:Didentity} and \eqref{tr:StieltjesD} must retain the
phrase ``the continuation from the stated initial domain.''  Literal
summation at an accidentally improved boundary point is a different
operation.  The example is also a useful negative control for software
that evaluates parameterized infinite sums after specialization.

\subsection{Hypotheses and constants to preserve during integration}

\paragraph{Moving and fixed indices.}
The polynomial identity for a word with a letter fixed at $u=-m$
contains no information about derivatives in the transverse variable
$u$.  It can be differentiated in its independently free letters.
To differentiate in $u$, first use the entire identity
\eqref{tr:genericgap}, and then specialize.  The functions
$\ell_{m,r}$ are the additional data introduced by that operation.
The finite-family bound in Proposition~\ref{tr:finitefamily} concerns
one marked index; it does not establish the same bound for several
independently marked indices.

\paragraph{Strict endpoints and Bernoulli signs.}
For a strict unit-lattice gap the finite power sum is
\[
 \sum_{z<x<y}x^m
   =\frac{B_{m+1}(y)-B_{m+1}(z+1)}{m+1}.
\]
Accordingly $B_1=-1/2$ and $B_1(1)=1/2$ occur in different positions.
At $m=0$, the missing endpoints account for the last term in
\[
 \HH(s,0,t)=\HH(s-1,t)-\HH(s,t-1)-\HH(s,t).
\]
The analogous term in \eqref{tr:commutator} must also remain after
transverse differentiation.  Its sign is fixed by the Hurwitz shift
identity, not by a choice of regularization.

\paragraph{Centered cutoff constants and the two zero-order zeta values.}
The mixed-tail proof uses the cutoff $0\leq n<N$ and the constant
term in the variable $N$.  The equality with the spectral value is
proved in \eqref{mt:cutoff-value}; it is not a general rule for every
cutoff coordinate.  Two different zero-order zeta values appear:
\[
                \zeta(0,1/2)=0,\qquad \zeta(0)=-\frac12.
\]
The first makes the continued sums of the subtracted centered even
polynomials vanish.  The second encodes the surviving rational boundary
term in \eqref{mt:boundary-constant}.  Dropping either distinction
changes the moment formula.  At $d=a+b-1$, the two oriented
$\gamma/2$ terms cancel the diagonal $-\gamma$; an isolated
$\zeta(1)$ is never introduced.

\paragraph{Domains, branches, and spanning statements.}
Endpoint parameters stay in the right half-plane.  The twists satisfy
the consecutive-sum inequalities in \eqref{zbg:tube}; removable zero
collisions do not remove the remaining nonzero-period hyperplanes.
The polylogarithmic primitives are first fixed in the right half-plane
and then continued as their complete anchored expressions, as in
\eqref{mt:L-polylog}.  The finite list in \eqref{mt:finite-span} is
a proved rational spanning list.  Its at-most-$\max(a,b)$ count is not
an assertion of arithmetic independence or minimality.  The ordinary
moment theorem also does not imply ordinary-zeta closure of its first
spectral derivative, which contains additional global logarithmic data.

\subsection{Reproducible verification and the limits of its evidence}

The analytic proofs above establish the unrestricted identities.
The companion scripts provide finite exact checks and independent
high-precision diagnostics of the most sensitive formulas.  The
recorded coverage is summarized below.

\begin{center}
\small
\begin{tabular}{@{}>{\raggedright\arraybackslash}p{1.40in}rr>{\raggedright\arraybackslash}p{3.43in}@{}}
\toprule
Family & Exact & Numerical & What is compared\\
\midrule
Geometric kernels and zero blocks
 & 590 & 0
 & Endpoint composition, deconcatenation, kernel recurrence, finite
   intervals, inverse intervals, gap coefficients, and insertion sums.\\[4pt]
Mixed-tail moments
 & 460 & 30
 & 400 boundary identities and 60 finite summation-by-parts identities;
   direct Hurwitz heads and centered tail expansions at two cutoffs.\\[4pt]
Transverse identities
 & 0 & 15
 & Finite Gamma/Hurwitz gaps, noninteger endpoint tails, Barnes values,
   negative-outer-order formulas, and the boundary discrepancy above.\\[4pt]
Normalized Mellin derivatives
 & 0 & 24
 & Subtracted logarithmic quadrature versus derivatives of explicit
   elementary Mellin transforms for two independent test functions.\\[4pt]
Moment generators
 & 0 & 3
 & Anchored quadrature, explicit polylogarithmic primitives, and a Taylor
   sum using 65 exact moment coefficients in each comparison.\\[4pt]
\midrule
Total & 1,050 & 72 & Finite algebra and numerical diagnostics have
                         different evidentiary roles.\\
\bottomrule
\end{tabular}
\end{center}

The 590 kernel checks use rational arithmetic only.  Positive
fourth-power exponential nodes and quarter-integral endpoints permit
exact tests even when $a-b$ is nonintegral.  The 460 moment checks are
finite identities over rational symbolic coefficient rings.  In
particular, the finite summation-by-parts tests keep the two tail
constants symbolic.

The transverse diagnostics use 65 decimal working digits and include
the discrepancy \eqref{audit:continued}--\eqref{audit:literal} as a
deliberate boundary control.  A separate 80-digit test of
\eqref{tr:Mjet} uses two quadratic-polynomial exponential functions,
one with nonreal coefficients and decay parameter.  For each function,
the centers $m=0,1,3$ and derivative orders $r=1,2,3,4$ compare the
subtracted logarithmic integral with a numerical derivative of its
known elementary Mellin transform.  All 24 cases satisfy the scripted
$10^{-50}$ tolerance; the largest observed discrepancy is about
$3.10\times10^{-79}$.  These functions test the normalization lemma
independently of the relative harmonic kernel.

The moment diagnostics use 95 working
digits, with orders through $(6,8)$ and moment indices through $M=7$.
They evaluate a direct finite Hurwitz-zeta head and a separately
assembled centered asymptotic tail at $(N,K)=(48,32)$ and $(64,36)$.
The largest observed discrepancy at the larger cutoff is about
$1.16\times10^{-63}$.  The three generator comparisons include a
nonreal parameter; their largest recorded Taylor-sum discrepancy is
about $1.08\times10^{-89}$.  The 65 Taylor coefficients are exact
inputs to those three numerical comparisons, not 65 additional
independent exact tests.

These residuals are observed floating-point discrepancies, not rigorous
error bounds or interval enclosures.  Working precision alone is not
an accuracy guarantee.  The recorded diagnostics cover the stated
examples; they do not constitute numerical quadrature of every
multidimensional integral in Theorem~\ref{tr:integral}.  Uniform
integrability, simultaneous branch cancellation, and continuation at
arbitrary orders are established by the proofs, not by a finite test
count.  The delivery includes neither a proof-assistant formalization
nor an external referee report.

\subsection{Status of the broader conjectures}

The contribution ledger records R7 as answered in the specified
arbitrary-depth zero-decoration family, R8 as answered for one marked
index and all of its nonpositive-center derivatives, and R1 as answered
for the full two-factor even-tail sector.  Extra harmonic factors,
several independently marked indices, and additional global spectral
derivative reductions remain separate questions.

The Gaussian $S_6$ and revised $S_8$ candidates retain their
conjectural status \cite{Repository,IncomingParity}.  No new functional
identity or exact period certificate for either candidate is supplied
here.  Existing nonmembership results concern their explicitly named
formal relation spaces and must not be read as period nonvanishing
theorems.  The next research questions in
Section~\ref{research:section} begin at these stated boundaries.

% End sections/05_audit_verification.tex

% Begin sections/06_research.tex
\section{Further exact research questions}
\label{research:section}

The results above change the starting point of the next investigation.
The all-zero-block generator and the single marked transverse derivative
need no further conjectural existence statement. The unresolved issues
concern additional factors, several marked positions, arithmetic
reduction, and analytic continuation beyond the domains proved here.
The questions below specify an object to compute and a criterion for
progress. They are proposals, not assertions that the suggested
reductions hold.

\subsection{Higher products and logarithmic moments}

\paragraph{Q1. Four and more even-order tails.}
For an even number $2d\geq4$ of even orders $a_j\geq2$, consider
\[
 F_{\boldsymbol a}(s)=\sum_{n\geq0}(n+1/2)^{-s}
                         \prod_{j=1}^{2d}\zeta(a_j,n+1).
\]
The centered expansion has exponents
$2d-\sum_j a_j-2k$, so the same subtraction argument makes all even
spectral integers regular. The open arithmetic problem is to evaluate
every $F_{\boldsymbol a}(-2M)$ in an explicit finite list of
multiple-zeta coordinates whose size is independent of $M$.
The first useful case is $(a_1,a_2,a_3,a_4)=(2,2,2,4)$.
Finite summation by parts and ordering the surviving indices give a
concrete route, but the two-factor Euler reduction does not by itself
remove every higher-depth coordinate. A satisfactory theorem should
give the finite coordinate map first and prove every subsequent
reduction. For an odd number of even tails, possible poles occur at
even spectral integers; that problem additionally requires a specified
Laurent finite part and is a different extension.

\paragraph{Q2. Which logarithmic remainders cancel?}
Corollary~\ref{mt:spectralderivatives} gives a convergent exact formula
for every $F_{a,b}^{(r)}(-2M)$, including its finite Hurwitz correction.
The remaining question is whether there are nontrivial, explicitly
specified rational linear combinations of these derivatives whose
logarithmic remainder sums cancel or reduce to endpoint data.
Begin with $r=1$, $(a,b)=(2,4)$ and $(2,2)$, and retain the
$-\log2/6$ term in \eqref{mt:derivative-example}.
One should seek an identity at the level of the summands or a
parameter-dependent kernel; an integer relation among a few decimal
values would only suggest a target. Independent differentiation of a
tail order, such as $\partial_a F_{a,b}$ before specializing $a$,
is another problem because it differentiates both the Hurwitz tail and
its asymptotic coefficients.

\Needspace{5\baselineskip}
\paragraph{Q3. A smallest proved span for each pair of orders.}
Corollary~\ref{mt:closure} supplies at most $\max(a,b)$ coordinates for
the entire sequence. Determine all reductions of that list obtainable
from stated algebraic identities, and give an algorithm producing a
smaller \emph{proved} spanning list. There are two distinct questions:
the rank of explicit rational coefficient vectors in a chosen formal
zeta-product alphabet, and arithmetic independence of the resulting
numbers. Only the first is a finite linear-algebra computation without
additional period theorems. A uniform description for $(2,2q)$ would
be a useful initial result.

\subsection{Several moving positions and primitive calculus}

\paragraph{Q4. Two independently marked positions.}
Construct a jointly continued gap formula for
\[
 \left.\partial_u^r\partial_v^q
       \HH_{a,b}(s,u,v,t)\right|_{u=-m,\,v=-n},
 \qquad m,n\geq0,\quad r,q\geq1,
\]
and identify a finite family of decorated lower-depth objects that
remains sufficient after arbitrary unmarked nonpositive deletions.
The adjacent case $m=n=0$, $r=q=1$ has a precise finite-gap guide.
For two surviving positive real lattice values $y>z$ in a finite
unit-lattice interval, put
\[
 P_r(y,z)=\ell_{0,r}(z+1)-\ell_{0,r}(y).
\]
The elementary symmetric-polynomial identity gives
\[
 \sum_{y>i>j>z}\log i\,\log j
                 =\frac12\bigl(P_1(y,z)^2-P_2(y,z)\bigr).
\]
Here $P_1$ is the negative sum of the gap logarithms and $P_2$ is
their squared-logarithm sum. This finite identity suggests the shape
of a Gamma--Stieltjes formula. A theorem at arbitrary endpoints must
still identify its decorated continuation with the prescribed
relative quotient. The separate-position and unequal-center cases
may require more than products of single-gap primitives.

\paragraph{Q5. Ordered versus symmetrized transverse data.}
Equation~\eqref{tr:insertionderivative} evaluates the sum of the three
possible insertion positions using depth-two data. Determine an exact
decomposition of several marked derivatives into their symmetrized
part and the remaining ordered part. At two marked positions,
stuffle products provide explicit relations but do not identify all
ordered differences. A useful result would describe a finite basis
of these ordered differences modulo the stated stuffle and deletion
relations, with the endpoint normalization retained throughout.

\paragraph{Q6. Endpoint differentiation and normalized antiderivatives.}
Combine the transverse calculus with differentiation in $a$ and $b$.
The identities
\[
 \partial_x\zeta(u,x)=-u\zeta(u+1,x),\qquad
 \mathcal B_m'(x)=\mathcal B_{m-1}(x)
\]
suggest a closed transport system connecting moving harmonic orders,
polygamma factors, and the balanced primitive ladder. Determine that
system for a depth-three word with one marked position, then integrate
it between specified endpoint values. The target is a formula whose
Gamma or Barnes constants are fixed by the same relative quotient,
rather than by separately chosen indefinite integrals. At $r\geq2$,
the analogous system involves generalized Stieltjes primitives and
should state their additive normalization explicitly.

\subsection{Rational shifts and continuation of the generators}

\paragraph{Q7. Common rational shifts and colored tails.}
For rational $c>1/2$, replace the moment family by
\[
 \sum_{n\geq0}(n+c-1/2)^{-s}
                 \zeta(a,n+c)\zeta(b,n+c).
\]
Its centered parity remains available, but the removed polynomial
terms now have continued values $\zeta(-2j,c-1/2)$, which generally
do not vanish. Establish the exact analogue of the finite Bernoulli
formula with all these counterterms. Then determine the appropriate
colored or cyclotomic multiple-zeta reduction for rational $c$.
Unequal shifts in the two tails should be treated separately, since
they no longer share the same centered expansion. Any differentiation
of rational multiplication formulas must keep the conductor powers
and their logarithmic derivatives.

\paragraph{Q8. Rational endpoint Gamma--Stieltjes identities.}
Specialize the finite endpoint formula
\eqref{tr:negativeformula} to $a,b\in\Q_{>0}$ and apply Hurwitz
distribution and character projections. For example,
\[
 \sum_{j=1}^{q}\zeta(s,j/q)=q^s\zeta(s)
\]
gives a completely specified hierarchy after differentiation at
$s=-m$. Determine which linear combinations of transverse weighted
tails reduce to derivatives of Dirichlet $L$-functions and logarithms
of Gamma or Barnes values, and which still require ordered coordinates.
The equality of the finite endpoint expressions provides an exact
starting point; no conjectural independence of Stieltjes constants is
needed for this transport step.

\paragraph{Q9. Singular hyperplanes beyond the common twist tube.}
Theorem~\ref{zbg:entire} proves joint holomorphy on $\Omega_k$.
The kernel identifies the only remaining candidate singular
hyperplanes as nonzero exponential periods of consecutive sums.
Determine their residues or local branching, first at depth two,
including further cancellations at specific nonzero periods.
For integral $a-b$, the function $z^{a-b}$ is single valued near
every nonzero node, so all finite node collisions are removable;
for $a-b=N\geq0$ the finite exponential sum gives the same conclusion.
For nonintegral endpoint differences,
a precise local continuation theorem would explain which parts of
the unit decoration polydisc can be enlarged and on which sheets.
It should treat the complete relative expression, since separately
continued Lerch summands can conceal cancellations.

\paragraph{Q10. Functional identities for the all-moment generator.}
The ODE and polylogarithmic representation in
Theorem~\ref{mt:generator} reduce each fixed pair of orders to
finitely many anchored primitives. Seek recurrences linking different
pairs $(a,b)$, reflection or continuation identities in $t$, and exact
special values at nonzero arguments where the polylogarithms have a
useful algebraic argument. The common radius $2\pi$ also makes the
nearest singularities at $\pm2\pi i$ accessible. Their local
expansions may give exact coefficient relations or full asymptotic
expansions of the moment sequence. Estimates alone are a secondary
goal here; the primary target is a new identity between the generators.

\Needspace{10\baselineskip}
\subsection{Integration with the existing conjectural program}

\paragraph{Q11. An analytic bridge to the Gaussian $S_6$ and $S_8$ candidates.}
The source manuscript distinguishes its formal Cayley or period
coordinates from proved evaluations. The new generators should be
compared with the existing Gaussian kernels only through an explicit
analytic transformation: a convergent integral identity, a justified
parameter deformation, or a proved functional equation. Determine
whether such a transformation sends a specific conjectural remainder
to a Gamma--Stieltjes gap or a mixed-tail coordinate evaluated here.
Until that bridge is proved, the present exact results supply neither
a proof nor a refutation of the $S_6$ or revised $S_8$ formulas.

\paragraph{Q12. Formalization of the exact algebra before the analytic layer.}
The divided-difference product identity, finite endpoint composition,
strict-gap binomial coefficients, Bernoulli deletion, and finite
summation-by-parts formulas are suitable first formalization targets.
Their statements should use exact rational or symbolic arithmetic and
retain the boundary terms. The analytic layer then needs explicitly
stated normalized-Mellin continuation, local normal convergence,
and Hurwitz asymptotics. The accompanying programs are reproducible
checks of finite instances; they do not replace those quantified
theorems. A useful integration milestone is a mechanically checked
finite algebraic core with the analytic lemmas listed as precise
dependencies.

\medskip
For an efficient next investigation, the most concrete starting points
are the four-tail example in Q1, the adjacent two-marked gap in Q4,
and the rational-shift extension in Q7. Each starts from a finite
identity or a convergent kernel already available here. They also test
three different potential obstructions: higher arithmetic depth,
nonlinear ordered primitives, and nonvanishing endpoint counterterms.

% End sections/06_research.tex

\begin{thebibliography}{99}
% Begin references.tex
% Bibliography entries only; article.tex supplies thebibliography.
% Repository links are immutable at the inspected commit.

\bibitem{Repository}
ProveIt Contributors,
\emph{Polylogarithms and their Arithmetic Bridges}.
ProveIt repository, pinned manuscript and editorial-ledger snapshot,
11 October 2026.
\url{https://github.com/VladimirReshetnikov/ProveIt/tree/7bd45777a8a127e5dc69aff8887ca36a79a3059d/Analysis/Polylogarithms/docs/manuscript}.

\bibitem{IncomingParity}
ProveIt research collection,
\emph{Harmonic Parity and Resolvent Identities:
Exact Lerch generators, Stieltjes primitives, and the quartic and higher
Tornheim layers}.
Research report, 11 October 2026; in particular, the centered harmonic
calculus, gap-deletion formulas, and research questions R1, R7, and R8.
Incoming archive at the revision of \cite{Repository}:
\url{https://github.com/VladimirReshetnikov/ProveIt/blob/7bd45777a8a127e5dc69aff8887ca36a79a3059d/docs/incoming/ProveIt_Harmonic_Parity_and_Resolvent_Identities_2026-10-11.zip}.

\bibitem{IncomingIndependent}
ProveIt research collection,
\emph{Independent Orders and Exact Reductions in Polylogarithm--Zeta
Calculus}.
Research report, 11 October 2026; in particular, the entire relative
Hurwitz interpolant and its nonpositive-index reductions.
Incoming archive at the revision of \cite{Repository}:
\url{https://github.com/VladimirReshetnikov/ProveIt/blob/7bd45777a8a127e5dc69aff8887ca36a79a3059d/docs/incoming/ProveIt_Independent_Orders_2026-10-11.zip}.

\bibitem{IncomingHigherReflection}
ProveIt research collection,
\emph{Higher Reflection Identities and Resonant Zeta Calculus}.
Research report, 11 October 2026; in particular, the centered harmonic
parity and resonant-continuation conventions.
Incoming archive at the revision of \cite{Repository}:
\url{https://github.com/VladimirReshetnikov/ProveIt/blob/7bd45777a8a127e5dc69aff8887ca36a79a3059d/docs/incoming/ProveIt_Higher_Reflection_Identities_2026-10-11.zip}.

\bibitem{INY}
Kentaro Ihara, Yayoi Nakamura, and Shuji Yamamoto,
\emph{Interpolant of truncated multiple zeta functions}.
The Ramanujan Journal \textbf{67} (2025), article 1.
\url{https://doi.org/10.1007/s11139-025-01045-2}.
Author preprint: \url{https://arxiv.org/abs/2407.20509}.
Especially Proposition 2.1, Proposition 3.3, Theorem 4.1,
and Corollary 4.2.

\bibitem{Bradley2007}
David M. Bradley,
\emph{A signed analog of Euler's reduction formula for the double zeta
function}.
Author preprint, arXiv:0707.4486 (2007).
\url{https://arxiv.org/abs/0707.4486}.
The unsigned specialization is the classical odd-weight Euler reduction;
see equations (2)--(3) and Proposition 1.

\bibitem{DLMFHurwitz}
National Institute of Standards and Technology,
\emph{NIST Digital Library of Mathematical Functions},
\S25.11, Hurwitz zeta function.
In particular, the shift law, polygamma and Bernoulli special values, parameter
differentiation, and Lerch's log-Gamma identity in
equations 25.11.3, 25.11.12, 25.11.14, 25.11.17, and 25.11.18;
the asymptotic expansion is 25.11.43.
Accessed 11 October 2026.
\url{https://dlmf.nist.gov/25.11}.

\bibitem{DLMFGamma}
National Institute of Standards and Technology,
\emph{NIST Digital Library of Mathematical Functions},
Chapter 5, Gamma function; especially \S5.2 (definitions),
\S5.7 (series expansions), \S5.11 (asymptotic expansions),
and \S5.15 (polygamma functions).
Accessed 11 October 2026.
\url{https://dlmf.nist.gov/5.2};
\url{https://dlmf.nist.gov/5.7};
\url{https://dlmf.nist.gov/5.11};
\url{https://dlmf.nist.gov/5.15}.

\bibitem{DLMFBernoulli}
National Institute of Standards and Technology,
\emph{NIST Digital Library of Mathematical Functions},
\S24.4, Basic properties of Bernoulli and Euler polynomials.
In particular, the difference, reflection, power-sum, and finite-expansion
identities in equations 24.4.1, 24.4.3, 24.4.9, and 24.4.12.
Accessed 11 October 2026.
\url{https://dlmf.nist.gov/24.4}.

\bibitem{DLMFPolylog}
National Institute of Standards and Technology,
\emph{NIST Digital Library of Mathematical Functions},
\S25.12, Polylogarithms.
Accessed 11 October 2026.
\url{https://dlmf.nist.gov/25.12}.

\bibitem{DLMFLerch}
National Institute of Standards and Technology,
\emph{NIST Digital Library of Mathematical Functions},
\S25.14, Lerch's transcendent.
Accessed 11 October 2026.
\url{https://dlmf.nist.gov/25.14}.

\bibitem{DLMFBarnes}
National Institute of Standards and Technology,
\emph{NIST Digital Library of Mathematical Functions},
\S5.17, Barnes' $G$-function (double gamma function).
In particular, the recurrence, normalization, product, and logarithmic
integral formulas in equations 5.17.1, 5.17.3, and 5.17.4.
Accessed 11 October 2026.
\url{https://dlmf.nist.gov/5.17}.

\bibitem{EspinosaMoll}
Olivier Espinosa and Victor H. Moll,
\emph{A generalized polygamma function}.
Integral Transforms and Special Functions \textbf{15} (2004), 101--115.
Author preprint: \url{https://arxiv.org/abs/math/0305079}.
The balanced negapolygamma normalization supplies the fixed constants
used in the primitive ladder.

\bibitem{Vardi}
Ilan Vardi,
\emph{Determinants of Laplacians and multiple gamma functions}.
SIAM Journal on Mathematical Analysis \textbf{19} (1988), 493--507.
\url{https://doi.org/10.1137/0519035}.
The Hurwitz-zeta derivative identity fixes the Barnes-$G$ normalization
used in the endpoint formula.

% End references.tex

\end{thebibliography}
\end{document}

